% concatenated from main.tex
\documentclass[reqno,11pt]{amsart}
\usepackage[top=2cm]{geometry}
\usepackage{amssymb}
\usepackage{amsthm}
\usepackage{yhmath}
\usepackage{amsfonts}
\usepackage{mathtools}
\usepackage{appendix}
\usepackage[dvipsnames]{xcolor}
\usepackage{enumitem}
\usepackage[colorlinks, 
linkcolor=blue,
anchorcolor=magenta,
citecolor=blue
]{hyperref}
\usepackage{orcidlink}
\newcommand{\supp}[1]{\mathrm{Supp}(#1)}
\newcommand{\e}{\mathbf{e}}

\newcommand{\norm}[2]{\Vert #1 \Vert_{#2}}

\DeclareMathOperator{\diag}{diag}
\DeclareMathOperator{\Div}{div}
\DeclareMathOperator{\sign}{sign}
\DeclareMathOperator{\curl}{curl}
% \usepackage{todonotes}
% \usepackage{showkeys}

%\usepackage[margin=1.5in]{geometry}
\usepackage[myheadings]{fullpage}

\theoremstyle{plain}
\newtheorem{theorem}{Theorem}[section]
\newtheorem{corollary}[theorem]{Corollary}
\newtheorem{lemma}[theorem]{Lemma}
\newtheorem{proposition}[theorem]{Proposition}
\theoremstyle{definition}
\newtheorem{definition}[theorem]{Definition}
\newtheorem{conjecture}[theorem]{Conjecture}
\newtheorem{example}[theorem]{Example}
\newtheorem{remark}[theorem]{Remark}
\newtheorem{case}{Case}[section]


\title[Inviscid shallow water equations]{Local Well-Posedness of the Inviscid Shallow Water Equations with Surface Tension}
\date{\today}
\author[]{Xin Liu \orcidlink{0000-0001-7021-8644}}
\address[Xin Liu]{
    Department of Mathematics, Texas A\$M University, College Station, TX, 77843, USA.
}
\email{xliu23@tamu.edu}

\author[]{Bingheng Yang$^*$ \orcidlink{0009-0005-0380-1596}}
\address[Bingheng Yang]{
	Department of Mathematics, Texas A\&M University, College Station, TX, 77843, USA.
}
\email{bhyang@tamu.edu}


\numberwithin{equation}{section}


\begin{document}


\begin{abstract}
We show that the inviscid shallow water equations with gravity and surface tension 
are locally well-posed in dimensions one and two. In contrast to previous works by Benzoni-Gavage, Danchin, and Descombes \cite{BenzoniGavageDanchinDescombes1D,BenzoniGavageDanchinDescombes3D}, we work only in the real-valued phase space and make use of the symmetry of the system. Our alternative proof for local well-posedness sheds some light on the symmetric structure of the shallow water system and can be used for further study, including designing numerical scheme, dispersion estimate, etc.

\smallskip 

{\noindent\bf Keyboards:} Euler-Korteweg, shallow water equations, surface tension, local well-posedness, symmetric systems.

{\noindent\bf MSC2020:} 35Q31, 35Q35, 76N10, 76B15, 76B45

\end{abstract}


\maketitle


{
  \hypersetup{linkcolor=blue}
  \tableofcontents
}

\section{Introduction}
We consider the system, %for $ n = 1,2 $,  
\begin{subequations}
\label{ShallowWaterEquations}
\begin{align}
\label{MaSs}
\partial_t h+\Div(hu)& = 0 \quad && \text{in} \quad \mathbb{R}^n, \\
\label{MoMeNtUm}
\partial_t u+u\cdot\nabla u+\nabla h-\nabla \Delta h & =0 \quad &&\text{in} \quad \mathbb{R}^n, \\
(h,u)\vert_{t=0} & =(h_0,u_0), 
\end{align}
\end{subequations}
in dimensions $ n = 1,2 $. We can view \eqref{ShallowWaterEquations} as a model for shallow water equations with surface tension. Here $h > 0$ denotes the height, $u \in \mathbb R^n$ denotes the velocity of the water wave, the pressure term  $\nabla h$ denotes the effect of gravity, and the third order term $-\nabla\Delta h$ arises from surface tension. Physically, \eqref{MaSs} and \eqref{MoMeNtUm} stand for the conservation of mass and momentum, respectively. The shallow water equations \eqref{ShallowWaterEquations} were formally derived from the free boundary problem for the incompressible water wave equations by the first author and his collaborators in \cite{LiLiuPeschka}, and were also previously known in the review article \cite{BreschDesjardinsMetivier} and the monograph \cite{Bresch}. As a shallow water model, the validity of \eqref{ShallowWaterEquations} was justified by Bresch and Noble in \cite{BreschNoble} in some scales. 

\smallskip
\begin{remark}
For $ n \geq 3 $, one should interpret \eqref{ShallowWaterEquations} as a special case of the Euler--Korteweg system \cite{BenzoniGavageDanchinDescombes3D, AntonelliMarcati1,BertiMasperoMurgante}, which models a fluid with quantum effects. See Section \ref{sec:irrotational-flow} for our remark about the irrotation flows in $ n $ dimension. 
\end{remark}
% as follows:
% \begin{subequations}
%     \label{EulerKorteweg}
%     \begin{gather}
%     \partial_t \rho+\Div (\rho u)=0, \\
%     \partial_t u+u\cdot \nabla u+\frac{1}{\rho}\nabla p(\rho)=\Div \mathbb K(\rho), \\
%     (\rho,u)|_{t=0}=(\rho_0,u_0),
%     \end{gather}
% \end{subequations}
% with $\rho>0$ being the density of the fluid, $p(\rho)=\frac{1}{2}\rho^2$ being the pressure and $\mathbb K(\rho)=(\Delta \rho) \mathbb {I}_{n\times n}$ being the so-called Korteweg stress tensor. 

\subsection{Main result}
Our goal is to prove local well-posedness for \eqref{ShallowWaterEquations}. We note that the cases of $ n = 1 $ and $ n = 2 $ admit fundamentally different structures. Therefore, we will treat them separately.   % (Theorem \ref{MainTheoremLWP}) for the system \eqref{ShallowWaterEquations}. 

In the case of $ n = 1 $, system \eqref{ShallowWaterEquations} is reduced to the following:
\begin{subequations}
    \label{sys:1-d_sw}
    \begin{align}
        \label{1-d-Mass}
        \partial_t h + \partial_x(hu) & = 0 \quad &&\text{in} \quad  \mathbb R, 
        \\
        \label{1-d-momentum}
        \partial_t u+ u \partial_x  u + \partial_x h - \partial_x^3 h & = 0 \quad && \text{in} \quad \mathbb R,\\
        \label{1-d-initial}
        (h,u)\vert_{t=0} & = (h_0,u_0). 
    \end{align}
\end{subequations}

\begin{theorem}
    \label{thm:1-d} Consider initial data $ (h_0, u_0)$ for system \eqref{sys:1-d_sw}, with
\begin{equation}
\label{h_initial-bound}
    0 < \underline h := \inf_{x} h_{0}(x) < 1 < \overline h := \sup_x h_{0}(x) < \infty,
\end{equation}
and
\begin{equation}
    \label{1-d-initial-regularity}
    h_0 - 1 \in H^3(\mathbb R), \quad  u_0 \in H^2(\mathbb R). 
\end{equation}
Then there exists $ T \in (0,\infty) $, depending only on the initial data, such that there exists a unique solution $ (h-1,u) \in L^\infty([0,T];H^3(\mathbb R)) \times L^\infty([0,T];H^2(\mathbb R)) $ to system \eqref{sys:1-d_sw}, with 
    \begin{equation}
        \label{1-d-a-priori-estimate}
        \sup_{0\leq t\leq T} \sum_{%\substack{%i+j\leq 2, \\ 0\leq i\leq 1, \\ 
        0\leq j\leq 2
        % }
        }\norm{\partial_x^j(h(t)-1),\partial_x^j u(t), \partial_x^{j+1} h(t) }{L^2} < C <\infty,
    \end{equation}
    for some constant $ C \in (0,\infty) $ depending only on the initial data. 
\end{theorem}

%Theorem \ref{thm:1-d} has the following higher dimensional analog.

Meanwhile, in the case of $ n=2 $, we will show
\begin{theorem}
    \label{MainTheoremLWP}
 Consider initial data $ (h_0, u_0) $ for system \eqref{ShallowWaterEquations}, with
\begin{equation}
    0 < \underline h := \inf_{x} h_{0}(x) < 1 < \overline h := \sup_x h_{0}(x) < \infty,
\end{equation}
and with regularity
\begin{equation}
    \label{n-d-initial-regularity}
    h_0 - 1 \in H^{7}(\mathbb R), \quad  u_0 \in H^{6}(\mathbb {R}^2). 
\end{equation}
Then there exists $ T \in (0,\infty) $ depending only initial data, such that there exists a unique solution $ (h-1,u) \in L^\infty([0,T];H^{4}(\mathbb R)) \times L^\infty([0,T];H^{3}(\mathbb {R}^2)) $ to system \eqref{ShallowWaterEquations}, with 
    \begin{equation}
        \sup_{0\leq t\leq T} \sum_{0\leq j\leq 3}\norm{\nabla^{j}(h(t)-1),\nabla^{j} u(t), \nabla^{j+1} h(t) }{L^2} < C <\infty,
    \end{equation}
    for some constant $ C \in (0,\infty) $ depending only on the initial data.
\end{theorem}

To prove Theorems \ref{thm:1-d} and \ref{MainTheoremLWP}, we will make use of the symmetric structure of system \eqref{ShallowWaterEquations}, which can be seen via the momentum variable $m:=hu.$ Indeed, one can write down the following equivalent system of the shallow water equations using $ (h,m) $: 
 % This is due to the underlying symmetric structure in \eqref{ShallowWaterEquations}: Namely, we make the substitution $m:=hu$ and apply $\nabla$ to \eqref{MaSs} to obtain the system
    \begin{subequations}
    \label{ShallowWaterMomentum}
    \begin{align}
    \label{ShallowWaterMomentumMass}
        \partial_t h+ \Div  m  & =0, \\
        \label{ShallowWaterMomentumMomentum}
        \frac{1}{h}\partial_t m +\frac{m}{h^2}\cdot\nabla m+ \frac{m}{h^2}\Div m+\nabla h-\nabla \Delta h & =\frac{(m\cdot\nabla)h}{h^3}  m, \\
        \label{ShallowWaterSurfaceGradient}
        \partial_t \nabla h+\nabla \Div m & =0.
    \end{align}
    \end{subequations}
% In these new coordinates, Theorem \ref{MainTheoremLWP} is equivalent to the following:
% \begin{theorem}
%     \label{MainTheoremLWPMomentum}
%        Let $s\in\mathbb{N}, \ s>1+\frac{n}{2}$ and consider the initial data $ (h_0, m_0) $ for system \eqref{ShallowWaterMomentum}, with
% \begin{equation}
%     0 < \underline h := \inf_{x} h_{0}(x) < 1 < \overline h := \sup_x h_{0}(x) < \infty,
% \end{equation}
% and
% \begin{equation}
%     h_0 - 1 \in H^{s+1}(\mathbb {R}), \quad  m_0 \in H^s(\mathbb {R}^n). 
% \end{equation}
% Then there exists $ T \in (0,\infty) $ depending only initial data, such that there exists a unique local-in-time solution $ (h-1,u) \in L^\infty([0,T];H^{s+1}(\mathbb R)) \times L^\infty([0,T];H^s(\mathbb {R}^n)) $ to system \eqref{sys:1-d_sw}, with 
%     \begin{equation}
%     \label{n-d-a-priori-estimate}
%         \sup_{0\leq t\leq T} \sum_{\substack{i+\vert \alpha \vert\leq \lceil \frac{n}{2}\rceil+1,\\ 0\leq i\leq \lceil \frac{n}{2}+\frac{1}{10}\rceil, \\ 0\leq \vert \alpha \vert\leq \lceil \frac{n}{2}\rceil +1}}\norm{\partial_t^i\partial^{\alpha}(h-1),\frac{1}{\sqrt h}\partial_t^i\partial^{\alpha} m, \partial_t^i\partial^{\alpha+1} h }{L^2} < C <\infty,
%     \end{equation}
%     for some constant $ C \in (0,\infty) $ depending only on the initial data. Here $\partial\in \{\partial_1,\ldots,\partial_n\}$ denotes a spatial derivative.
% \end{theorem}
% Let us motivate Theorem \ref{MainTheoremLWPMomentum}.

\smallskip 

To explain the idea, in the case of $ n = 1 $, i.e., in one spatial dimension, let $ U:= (h-1, m, \partial_x h)^\top \in \mathbb R^3 $. One can write system \eqref{sys:1-d_sw} (or equivalently, system \eqref{ShallowWaterMomentum}) as
% Note that in dimension one, \eqref{ShallowWaterMomentum} is a \emph{symmetric} system:
\begin{subequations}
\label{ShallowWaterMomentum1DMatrix}
    \begin{gather}
    \begin{pmatrix}
        1 & 0 & 0 \\
        0 & \frac{1}{h} & 0 \\
        0 & 0 & 1 \\
    \end{pmatrix}\partial_t U+\begin{pmatrix}
        0 & 1 & 0 \\
        1 & \frac{2m}{h^2} & 0 \\
        0 & 0 & 0 \\
    \end{pmatrix}\partial_x U+\begin{pmatrix}
        0 & 0 & 0 \\
        0 & 0 & -1 \\
        0 & 1 & 0 \\
    \end{pmatrix}\partial_x^2 U=\begin{pmatrix}
        0 \\
        \frac{m^2}{h^3}\partial_x h \\
        0 \\
    \end{pmatrix}, \\
    U(0,x)=(h_0-1,m_0,\partial_x h_0).
\end{gather}
\end{subequations}
One can see that system \eqref{ShallowWaterMomentum1DMatrix} resembles a symmetric hyperbolic system, with additional second order terms. In particular, the associated linear operator generates a bounded group, and therefore there is no loss of regularity; see Lemma \ref{LemmaPhi}, below. Due to this symmetry, the local well-posedness of the generalized system \eqref{ShallowWaterV} follows easily from the method of vanishing viscosity. We refer to Section \ref{vv} for the complete argument. At last, in Section \ref{Sect:WP}, we show the well-posedness of system \eqref{ShallowWaterMomentum1DMatrix} (or equivalently, system \eqref{sys:1-d_sw}) using an approximating argument.

\smallskip 

However, in dimension two, the symmetric structure of the shallow water system \eqref{ShallowWaterEquations} differs from that in \eqref{ShallowWaterMomentum1DMatrix}. Indeed, only the acoustic component of the flow admits a symmetric structure, while the vorticity component of the flow is transported. See Section \ref{sec:2d} for the detailed formulation of $ n = 2 $, as well as the \emph{a priori} estimates for the 2d system \eqref{ShallowWaterMomentum2D-1}. 
In Section \ref{Regularity-remark} we explain why the approximating argument in Section \ref{Sect:WP} cannot be directly applied to the 2d case without additional regularity assumptions.
Finally, in Section \ref{sec:irrotational-flow}, we remark about a result for irrotational flows in dimensions $\geq 3$. 

% { {\bf To check this part later} \color{purple} 
% On the other hand, the original system \eqref{ShallowWaterEquations} is not. We henceforth exploit the symmetric structure of \eqref{ShallowWaterMomentum1DMatrix}, and prove the equivalent Theorem \ref{MainTheoremLWPMomentum} instead. The proof of Theorem \ref{MainTheoremLWPMomentum} is based on a suitable energy estimate, and is then followed by a vanishing viscosity argument. The key difficulty is to obtain energy inequalities as in Proposition \ref{EnergyProposition} and in Proposition \ref{EnergyPropositionNd}: In dimension one, all matrix coefficients in \eqref{ShallowWaterMomentum} are symmetric or skew-symmetric, and furthermore carry favorable signs. As a result, there are no loss of derivatives in our energy estimates and the proof of Theorem \ref{MainTheoremLWPMomentum} turns out to be, more or less, straightforward. However, we lose the symmetric structure in higher dimensions due to the contribution from \eqref{BadTerm}. To deal with this, we symmetrize and introduce the good variable \eqref{Curled} that satisfies a simple transport equation \eqref{CurledTransport}, which provides $H^s$ estimates and thus helps us to control the skew-symmetric contributions. Finally, in dimensions $\geq 4$, we define the $\curl$ operator appropriately so that the same analysis can be carried out in higher dimensions. 
% }



\subsection{Background and literature on inviscid water waves with gravity and surface tension} %The role of the dynamical contact angles/lines has been extensively studied by the physical science community. We refer the reader to the following physics literature \cite{Blake,BlakeRuschak,Bertozzi,DeGennes,DiezKondicBertozzi,Dussan,DussanRameGaroff} for theoretical, numerical and experimental studies, as well as applications to wetting theory. 
For inviscid shallow water equations \emph{without} surface tension, or for \emph{viscous} shallow water equations with surface tension, we refer the reader to Bresch's book \cite{Bresch}, Lannes' works \cite{LannesBook,Lannes} and the review article by Bresch--Desjardins--M{\'e}tivier \cite{BreschDesjardinsMetivier} for an extensive overview on the mathematical properties of shallow water models in the absence/presence of boundary. In particular, Bresch and Noble in \cite{BreschNoble} rigorously justified the validity of \eqref{ShallowWaterEquations} in smooth function spaces.

% Curiously, although Bresch and Noble in \cite{BreschNoble} were able to rigorously justify the validity of \eqref{ShallowWaterEquations} in smooth function spaces, we were unable to find pre-existing well-posedness results for \eqref{ShallowWaterEquations} in the context of inviscid shallow water models with surface tension. \\

Much more mathematical studies on inviscid water waves are available without the assumption of shallowness. Here we highlight the following works: Wu \cite{WUI,WUII,WUIII,WUIV}, Coutand--Shkoller \cite{CoutandShkoller}, Alazard--Burq--Zuily \cite{AlazardBurqZuily1,AlazardBurqZuily2,AlazardBurqZuily3} and Deng--Ionescu--Pausader--Pusateri \cite{IonescuPusateriI,IonescuPusateriII,IonescuPusateriIII,DengIonescuPausaderPusateri} managed to show local and global well-posedness of incompressible irrotational water waves with or without surface tension using various methods, whereas Castro--C{\'o}rdoba--Fefferman--Gancedo--G{\'o}mez-Serrano \cite{Castro1,Castro2} investigated the finite-time formation of splash singularities. Finally, we note that Ambrose--Masmoudi \cite{AmbroseMasmoudi1,AmbroseMasmoudi2} and Agrawal \cite{Agrawal} investigated the zero surface tension limit problem. %However, we remark here that unlike any of the aforementioned works, our model \eqref{ShallowWaterEquations} is compressible and we allow our flow to have nonzero vorticity in higher dimensions. As a result, much of the results mentioned above that involve dispersive estimates cannot be directly translated into our framework. In fact, we see that the curl \eqref{Curled} is a good unknown for us.\\

%Regarding inviscid quantum fluid models, we refer the reader to the works of Bresch--Gisclon--Lacroix-Violet \cite{BreschGisclonLacroixViolet} and Antonelli--Marcati \cite{AntonelliMarcati2,AntonelliMarcati3} for a comprehensive outlook regarding the theory of Euler--Korteweg system, as well as its relation to capillary effects and quantum hydrodynamics. In this context, Benzoni-Gavage--Danchin--Descombes \cite{BenzoniGavageDanchinDescombes1D,BenzoniGavageDanchinDescombes3D} obtained strong local well-posedness of the Euler--Korteweg systems with a general class of Korteweg stress tensors $\mathbb K$, and Audiard--Haspot \cite{AudiardHaspot} showed global well-posedness in dimensions 3 or above for irrotation flow with small data. In a related work, Benzoni-Gavage and Chiron \cite{BenzoniGavageChiron} investigated the asymptotic behavior of the Euler--Korteweg system. Finally, we mention the works by Gwiazda et al. \cite{DebiecGwiazdaGwiazdaTzavaras,GallenmullerGwiazdaGwiazdaWoznicki} and Giesselmann--Tzavaras \cite{GiesselmannTzavaras} on the weak solutions of the Euler--Korteweg system and relations to the Cahn--Hilliard equation.

% \begin{remark}
    We remark here that Theorem \ref{thm:1-d} and Theorem \ref{MainTheoremLWP} corresponds to the $K=1$ case for the Euler--Korteweg system in \cite[Theorem 1.1]{BenzoniGavageDanchinDescombes3D}. The proof in \cite{BenzoniGavageDanchinDescombes3D} relies on finding the good variable 
    \[
    w:=\frac{\nabla h}{\sqrt{h}},
    \]
    then complexifying the equations and performing rather involved weighted energy estimates. Moreover, the argument in \cite{BenzoniGavageDanchinDescombes3D} requires a vanishing viscosity argument using the hyperdissipative term $-\varepsilon \Delta^2$, for which uniform-in-$\varepsilon$ estimates are more difficult to establish. We provide a shorter and more physically intuitive proof based on another good variable \eqref{ThetaVariable}, which in turn gives more insights to the hidden symmetric structure of \eqref{ShallowWaterEquations}. In addition, we believe that our approach based on (almost) symmetric systems can be useful for further numerical studies of shallow water waves. See for instance \cite{BreschEtAl,EiterGiesselmannLasarzikOffnerSauerborn,Guermond1, Guermond2,Guermond3,NobleVila,Xing,XingShu,XingZhangShu} for numerical studies of shallow water models and of the Euler--Korteweg system.
% \end{remark}


\subsection{Notations}
\label{Notations}
In this paper, we denote $A\lesssim B $ if
\[
A\leq CB,
\]
for some generic constant $C \in (0,\infty)$ that may be different from line to line. We will use subscripts for $ C $ to emphasize the dependency. \\

In addition, for $f_1,\cdots,f_n\in X$, we denote
\[
\Vert f_1,\cdots,f_n\Vert_{X}^2:=\sum_{j=1}^n\Vert f_j\Vert_{X}^2.
\]
Here $ X $ is some Banach space. We denote
\[
\int (\cdot)d\vec{x}:=\int_{\mathbb R^n}(\cdot) d\vec{x}=\int_{\mathbb R^n} (\cdot) dx_1\ldots dx_n,
\]
for integrals on $\mathbb R^n$, and the pairing
\[
(\cdot,\cdot)
\]
stands for the standard $L^2$ inner product. We will use
\begin{equation}
    \label{partial-z}
    \partial \in \{\partial_{x_1},\partial_{x_2},\cdots,\partial_{x_n}\},
\end{equation}
to denote a spatial derivative along arbitrary direction. Lastly, the notation
\[
(\cdot)^\top
\]
stands for the transpose of a matrix.


\section{\textit{A priori} estimates in 1D}
\label{1DEnergyEstimates}
In this section, we focus on the {\it a priori} energy estimates for the 1d shallow water system. We write \eqref{ShallowWaterMomentum1DMatrix} as
\begin{equation}
\label{ShallowWaterMomentumConsice}
A_0(U)\partial_t U+A_1(U)\partial_x U+A_2\partial_x^2 U=S(U),
\end{equation}
where
\begin{equation}
    \label{def:U}
    U := \begin{pmatrix}
        h - 1 \\ m \\ \partial_x h
    \end{pmatrix},
\end{equation}
and
\begin{equation}\label{def:As}
\begin{gathered}
    A_0(U) := \begin{pmatrix}
        1 & 0 & 0 \\
        0 & \frac{1}{h} & 0 \\
        0 & 0 & 1 \\
    \end{pmatrix}, \quad A_1(U):=\begin{pmatrix}
        0 & 1 & 0 \\
        1 & \frac{2m}{h^2} & 0 \\
        0 & 0 & 0 \\
    \end{pmatrix}, \quad A_2:=\begin{pmatrix}
        0 & 0 & 0 \\
        0 & 0 & -1 \\
        0 & 1 & 0 \\
    \end{pmatrix},  \\
    S(U):=\begin{pmatrix}
        0 \\
        \frac{m^2}{h^3}\partial_x h \\
        0 \\
    \end{pmatrix}. 
    \end{gathered}
\end{equation} 
We will make use of the symmetric structure of equation \eqref{ShallowWaterMomentumConsice}. 



\subsection{\textit{A priori} assumption,  the generalized system, and the initial data}
\label{subsec:generalized-system-initial}

Instead of \eqref{ShallowWaterMomentumConsice},
it is more convenient to consider a generalized system.
Let
\begin{equation}
\label{Def:V}
\begin{gathered}
V:=\begin{pmatrix}
    \eta -1 \\
    \omega \\
    \phi
\end{pmatrix}, \quad A_0=A_0(V):=\begin{pmatrix}
    1 & 0 & 0 \\
    0 & \frac{1}{\eta} & 0 \\
    0 & 0 & 1
\end{pmatrix}, \quad A_1=A_1(V):=\begin{pmatrix}
    0 & 1 & 0 \\
    1 & \frac{2\omega}{\eta^2} & 0 \\
    0 & 0 & 0
\end{pmatrix},  \\
{S}_{\phi} = S_\phi(V):=\begin{pmatrix}
    0 \\
    \frac{m^2}{\eta^3}\phi \\
    0 \\
\end{pmatrix},
\end{gathered}
\end{equation}
and the generalized system
\begin{equation}
    \label{ShallowWaterV}
    A_0\partial_t V+A_1\partial_x V+A_2\partial_x^2 V={S}_{\phi}.
\end{equation}
Notice that, as long as the solutions to \eqref{ShallowWaterMomentumConsice} and \eqref{ShallowWaterV} exist with the same initial data and are sufficient regular, it is easy to prove $ U(x,t) = V(x,t) $ for all $ 0 \leq t \leq T $ and all $x\in\mathbb{R}$, where $ T $ is the maximal existing time.  

For the generalized system \eqref{ShallowWaterV}, consider the energy functional
\begin{equation}
    \label{Energy}
    \mathcal E(t):= \sum_{\substack{i+j\leq 2,\\ i,j \geq 0}} \Vert\partial_t^i\partial_x ^j (\eta-1),{\partial_t^i\partial_x^j \omega},\partial_t^i\partial_x^{j}\phi\Vert_{L^2}^2,
\end{equation}
and the intermediate energy
\begin{equation}
    \label{intrmd-energy}
    \mathfrak E(t): = \sum_{\substack{i+j\leq 2,\\ i,j \geq 0}} (A_0(V) \partial_t^i \partial_x^j V (t), \partial_t^i \partial_x^j V(t)). 
\end{equation}
In particular, provided that 
\begin{equation}
    \label{non-degenerate}
    0<\frac{1}{2}\underline h \leq \inf_{x,t} \eta(x,t) \leq \sup_{x,t} \eta(x,t) \leq 2 \overline h, 
\end{equation}
there exists a constant $ C_{\underline h, \overline h} \in (0,\infty) $ depending on $ \underline h, \overline h $, such that
\begin{equation}
    \label{energy-eqvlt}
    \frac{1}{C_{\underline h, \overline h}}\mathcal E(t) \leq \mathfrak E(t) \leq C_{\underline h, \overline h} \mathcal E(t).
\end{equation}
Moreover, at $ t = 0 $, we choose initial data with regularity
\begin{equation}\label{V-initialdata} (\eta-1,\omega, \phi)\vert_{t=0} := (\eta_0 -1, \omega_0, \phi_0 ) \in H^3(\mathbb{R}) \times H^4(\mathbb{R}) \times H^4(\mathbb{R}), \end{equation} 
as the initial data for $ \partial_t \eta, \partial_t \omega, \partial_t \phi, \partial_{tt} \eta, \partial_{tt} \omega, \partial_{tt} \phi $ are explicitly given by
\begin{equation}
\label{initial_data_time}
\begin{aligned}
    \partial_t \eta(t=0) & = - \partial_x \omega_0,\\
    \partial_t \omega(t=0) &= - \frac{2\omega_0\partial_x \omega_0}{\eta_0}  - \eta_0 \partial_x \eta_0 + \eta_0 \partial_{x}^2 \phi_0 + \frac{\omega_0^2 \partial_x \eta_0}{\eta_0^2}, \\ 
    \partial_t \phi (t=0) & = - \partial_x^2 \omega_0, \\
    \partial_t^2 \eta(t=0) &= -2\frac{\omega_0 }{\eta_0^2} \phi_0\partial_x \omega_0-\frac{\omega_0^2}{\eta_0^2 }\partial_x \phi_0+\frac{2\omega_0^2}{\eta_0^3}\phi_0\partial_x\eta_0-\partial_x\eta_0 \partial_x^2\phi_0-\eta_0\partial_x^3 \phi_0+(\partial_x\eta_0)^2  \\
    & +\eta_0\partial_x^2 \eta_0+\frac{2(\partial_x\omega_0)^2}{\eta_0}+\frac{2\omega_0}{\eta_0}\partial_x^2 \omega_0-\frac{2\omega_0}{\eta_0^2 } \partial_x \eta_0\partial_x \omega_0, \\
     \partial_t^2 \phi (t=0) &= \partial_x \partial_t^2 \eta(t=0) \\
     & = -\frac{2\phi_0}{\eta_0^2 } (\partial_x \omega_0)^2 -\frac{4\omega_0}{\eta_0^2 } \partial_x \phi_0 \partial_x \omega_0-2\frac{\omega_0}{\eta_0^2 } \phi_0\partial_x^2 \omega_0+4\frac{\omega_0}{\eta_0^3} \phi_0\partial_x \eta_0 \partial_x \omega_0 -\frac{\omega_0^2}{\eta_0^2}\partial_x^2 \phi_0\\
     &+\frac{4\omega_0^2}{\eta_0^3} \partial_x \eta_0\partial_x \phi_0+\frac{4\omega_0}{\eta_0^3}\phi_0\partial_x \omega_0\partial_x \eta_0+\frac{2\omega_0^2}{\eta_0^3}\phi \partial_x^2 \eta_0 -\frac{6\omega_0^2}{\eta_0^4}\phi_0 (\partial_x \eta_0)^2-\partial_x^2\eta_0 \partial_x^2 \phi_0\\
     & -2\partial_x \eta_0 \partial_x^3 \phi_0-\eta_0\partial_x^4 \phi_0+3\partial_x\eta_0\partial_x^2\eta_0+\eta_0\partial_x^3 \eta_0+\frac{6}{\eta_0}\partial_x\omega_0\partial_x^2\omega_0-\frac{2(\partial_x\omega_0)^2}{\eta_0^2}\partial_x\eta_0+\frac{2\omega_0}{\eta_0}\partial_x^3\omega_0 \\
     & -\frac{2\omega_0}{\eta_0^2}\partial_x^2 \omega_0-\frac{2}{\eta_0^2}\partial_x\eta_0 (\partial_x\omega_0)^2-\frac{2\omega_0}{\eta_0^2}\partial_x^2\eta_0\partial_x \omega_0-\frac{2\omega_0}{\eta_0^2}\partial_x\eta_0\partial_x^2 \omega_0+\frac{4\omega_0}{\eta_0^3}(\partial_x \eta_0)^2 \partial_x \omega_0,\\
     \partial_t^2 \omega (t=0) &=-\frac{2\partial_x \omega_0}{\eta_0}\left(- \frac{2\omega_0\partial_x \omega_0}{\eta_0}  - \eta_0 \partial_x \eta_0 + \eta_0 \partial_{x}^2 \phi_0 + \frac{\omega_0^2 \partial_x \eta_0}{\eta_0^2}\right)\\
     &-\frac{2\omega_0}{\eta_0}\partial_x \left(- \frac{2\omega_0\partial_x \omega_0}{\eta_0}  - \eta_0 \partial_x \eta_0 + \eta_0 \partial_{x}^2 \phi_0 + \frac{\omega_0^2 \partial_x \eta_0}{\eta_0^2}\right)\\
     &-\frac{2\omega_0\partial_x \omega_0}{\eta_0^2}\partial_x\omega_0+\partial_x\eta_0\partial_x \omega_0 -\partial_x\omega_0\partial_x^2\phi_0-\eta_0\partial_x^4 \omega_0  \\
     & +\frac{2\omega_0\partial_x\eta_0}{\eta_0^2} \left(- \frac{2\omega_0\partial_x \omega_0}{\eta_0}  - \eta_0 \partial_x \eta_0 + \eta_0 \partial_{x}^2 \phi_0 + \frac{\omega_0^2 \partial_x \eta_0}{\eta_0^2}\right)-\frac{\omega_0^2}{\eta_0^2} \partial_x^2 \omega_0+\frac{2\omega_0^2\partial_x\eta}{\eta_0^3} \partial_x \omega_0.
\end{aligned}
\end{equation}
It follows from \eqref{initial_data_time} that 
$ \mathcal E(0) < \infty $ if \eqref{non-degenerate} and \eqref{V-initialdata} hold.% $ \eta_0 - 1 \in H^3(\mathbb{R})$ and $ \omega_0, \phi_0 \in H^4(\mathbb{R}) $. 

\subsection{\textit{A priori} estimates of the generalized system \eqref{ShallowWaterV}} 
We focus first on obtaining the {\it a priori} energy estimates for the solution $ V $ to system \eqref{ShallowWaterV}. We assume the non-degenerate condition \eqref{non-degenerate} holds
% following {\it a priori} bounds: 
% \begin{align}
%     \label{APrioriHeight}
%     \begin{aligned}
%     0<\frac{1}{2} \underline h \leq \eta(t,x)\leq 2 \overline h < \infty , \\
%     \end{aligned}
% \end{align}
for all $0\leq t \leq T$, where $T$ is the maximum existence time. We will justify \eqref{non-degenerate} in Section \ref{Justification} below.

Our goal of this section is to establish the following proposition:
\begin{proposition}
    \label{prop:a-priori-V}
    Under the assumption \eqref{non-degenerate}, the solution to \eqref{ShallowWaterV} satisfies
    \begin{equation}
        \label{est:a-priori-energy-0}
        \sup_{0\leq t \leq T} \mathfrak E(t) \leq 2 \mathfrak  E(0),
    \end{equation}
    and thus, thanks to \eqref{energy-eqvlt}, 
    \begin{equation}
        \label{est:a-priori-energy}
        \sup_{0\leq t \leq T} \mathcal E(t) \leq C_{\underline h, \overline h} \mathcal  E(0),
    \end{equation}
    for some $ T \in (0,\infty) $ and some $ C_{\underline h, \overline h} \in (0,\infty) $ depending only on the initial data and $ \underline h, \overline h $. 
\end{proposition}

The proof of Proposition \ref{prop:a-priori-V} relies on the following property of the linear operator $ A_0 \partial_t + A_1 \partial_x + A_2 \partial_x^2$:
% \subsection{Sobolev embedding in 1D}
% Recall the following form of Sobolev embedding in 1D:
% \begin{proposition}
%     Let $f\in H^s(\mathbb{R})$ with $s>\frac{3}{2}$. Then in fact $f\in C^1(\mathbb{R})$, and 
%     \[
%     \Vertf\Vert_{L^{\infty}(\mathbb{R})}\lesssim \Vertf\Vert_{H^1(\mathbb{R})}, \quad \Vertf\Vert_{C^1(\mathbb{R})}\lesssim \Vertf\Vert_{H^2(\mathbb{R})}.
%     \]
% \end{proposition}
% We also require the following lemma to close the energy estimates:
\begin{lemma}
\label{LemmaPhi}
    Let $\Phi: \mathbb R \times \mathbb R \to \mathbb R^3$ be the solution to the linear equation
    \begin{equation}
    \label{PhiEq}
    A_0\partial_t \Phi+A_1\partial_x \Phi+A_2\partial_x^2 \Phi=G,
    \end{equation}
    where $G\in L_t^{\infty}L^2_x$. Then for a.e. $ t $, we have
    \begin{equation}
    \label{LemmaPhiEstimate}
        \frac{d}{dt} (A_0 \Phi,\Phi) \leq C_{\underline h, \overline h} (\mathcal E(t) + 1) \Vert\Phi\Vert_{L^2}^2+ 2 \Vert G\Vert_{L^2} \Vert\Phi\Vert_{L^2}.
    \end{equation}
\end{lemma}
\begin{proof}
    Taking the $ L^2 $-inner product of \eqref{PhiEq} with $ 2 \Phi $ and integrating the resultant by part lead to 
    \begin{equation}
        \label{phi:001}
        \dfrac{d}{dt} (A_0 \Phi, \Phi) =  (\partial_t A_0 \Phi, \Phi) + (\partial_x A_1 \Phi, \Phi) + \underbrace{2(A_2 \partial_x \Phi, \partial_x \Phi)}_{=0} + 2 (G,\Phi),
    \end{equation}
    where we used the fact that $ A_0, A_1 $ are symmetric matrices and $ A_2 $ is a constant skew-symmetric matrix. Applying the Hölder inequality to the right-hand side of \eqref{phi:001} yields
    \begin{equation}
    \label{phi:002}
        \dfrac{d}{dt}(A_0 \Phi, \Phi) \leq \norm{\partial_t A_0, \partial_x A_1}{L^\infty}\norm{\Phi}{L^2}^2 + 2 \norm{G}{L^2} \norm{\Phi}{L^2}. 
    \end{equation}
    It remains to estimate $ \norm{\partial_t A_0, \partial_x A_1}{L^\infty} $. From \eqref{def:As} and \eqref{ShallowWaterV}, one has
    \begin{equation}
        \label{phi:003}
        \begin{gathered}
        \norm{\partial_t A_0, \partial_x A_1}{L^\infty} \leq C_{\underline h} \norm{\partial_t \eta, \partial_x \omega, \omega \partial_x \eta}{L^\infty} \leq C_{\underline h} (\norm{\partial_x \omega, \omega,\partial_x \eta}{L^\infty}^2 + 1) \\
        \leq C_{\underline h} (\sum_{j=0}^2 \norm{\partial_x^j \omega, \partial_x^j (\eta-1)}{L^2}^2 + 1) \leq C_{\underline h, \overline h} (\mathcal E(t) + 1),
        \end{gathered}
    \end{equation}
    where we have used the Sobolev embedding inequality $ \norm{\cdot}{L^\infty} \lesssim \norm{\cdot }{H^1} $. 
    This finishes the proof. 
\end{proof}
Now we are ready to establish the {\it a priori} estimates in the generalized system \eqref{ShallowWaterV}. Applying \eqref{LemmaPhiEstimate} to \eqref{ShallowWaterV} immediately leads to 
\begin{equation}
    \label{basic:energy}
    \dfrac{d}{dt}(A_0V,V) \leq C_{\underline h,\overline h} (\mathcal E(t) + 1) \norm{V}{L^2}^2 + 2 \norm{S_\phi}{L^2} \norm{V}{L^2}.
\end{equation}
Meanwhile, one has the following systems for $ V_{tt} := \partial_{t}^2 V, V_{xt} := \partial_x\partial_{t} V $, and $ V_{xx} := \partial_{x}^2 V $:
\begin{equation}
    \label{sys:dd-V}
    A_0 \partial_t V_{a} + A_1 \partial_x V_{a} + A_2 \partial_{x}^2  V_a = G_a, \qquad \text{for} ~ a \in \lbrace tt, xt, xx\rbrace,
\end{equation}
where
\begin{align}
    \label{id:tt-G}
    & \begin{aligned}
     & G_{tt} := \partial_{t}^2 S_\phi - 2 \partial_t A_0 \partial_{t}^2 V - \partial_{t}^2 A_0 \partial_t V - 2 \partial_t A_1 \partial_t\partial_{x}V- \partial_{t}^2 A_1 \partial_x V,
    \end{aligned}
    \\
    & \begin{aligned}
        & G_{xt} := \partial_{xt} S_\phi - \partial_x A_0 \partial_{tt} V - \partial_t A_0 \partial_{xt} V - \partial_{xt} A_0 \partial_t V \\
        & \qquad - \partial_x A_1 \partial_{xt}V - \partial_{t} A_1 \partial_{xx} V - \partial_{xt} A_1 \partial_x V, 
    \end{aligned}
    \label{id:xt_G}
    \\
    \label{id:xx_G}
    & \begin{aligned}
        & G_{xx} := \partial_{xx} S_\phi - 2 \partial_x A_0 \partial_{xt} V - \partial_{xx}A_0 \partial_t V - 2 \partial_x A_1 \partial_{xx} V - \partial_{xx} A_1 \partial_x V. 
    \end{aligned}
\end{align}
Then applying \eqref{LemmaPhiEstimate} in Lemma \ref{LemmaPhi} to \eqref{sys:dd-V}, one has that
\begin{equation}
    \label{h2:001}
    \dfrac{d}{dt} \sum_{a \in \lbrace tt, xt, xx\rbrace} (A_0 V_a, V_a) \leq C_{\underline h, \overline{\eta}} (\mathcal E(t) + 1) \norm{V_{tt}, V_{xt}, V_{xx}}{L^2}^2 + 2\norm{G_{tt},G_{xt},G_{xx}}{L^2} \norm{V_{tt}, V_{xt}, V_{xx}}{L^2}. 
\end{equation}
Moreover, one can directly calculate that 
\begin{equation}
\label{h-1-eq}
\begin{gathered}
    \dfrac{d}{dt}\lbrace (A_0 \partial_x V, \partial_x V) + (A_0 \partial_t V, \partial_t V) \rbrace = 2 (A_0 \partial_{x}\partial_t V, \partial_x V) + 2 (A_0 \partial_{t}^2 V, \partial_t V) \\ + (\partial_t A_0 \partial_x V, \partial_x V) + (\partial_t A_0 \partial_t V,\partial_t V)\\
    \lesssim \norm{A_0,\partial_t A_0}{L^\infty} \norm{\partial_x V, \partial_{x}\partial_tV,\partial_{t}V,\partial_{t}^2 V}{L^2}^2 \\
    \overset{\eqref{phi:003}}{\leq} C_{\underline h, \overline h} (\mathcal E(t) + 1)\norm{\partial_x V, \partial_{x}\partial_tV,\partial_{t}V,\partial_{t}^2 V}{L^2}^2.
    \end{gathered}
\end{equation}
Combining \eqref{basic:energy}, \eqref{h2:001}, and \eqref{h-1-eq} leads to
\begin{equation}
    \label{h-002}
    \dfrac{d}{dt}\mathfrak E(t)\leq C_{\underline h, \overline h} (\mathcal E(t) + 1)^2 + 2 \norm{S_\phi,G_{tt},G_{xt},G_{xx}}{L^2} (\mathcal E(t))^{\frac{1}{2}}.
\end{equation}
It remains to estimate $ \norm{S_\phi,G_{tt},G_{xt},G_{xx}}{L^2} $. However, this is immediate. Since $ A_0, A_1, S_\phi $ depend only on $ V $ but not the derivative(s) of $ V $, one can write that, for $ \Psi \in \lbrace A_0, A_1, S_\phi \rbrace $, 
\begin{align}
    \label{est:nonlinearity-001}
    & \norm{\partial_x \Psi}{L^\infty} \leq (\norm{\Psi}{L^{\infty}}+\norm{ \partial_x \Psi }{L^\infty}) \norm{\partial_x V}{L^\infty} \leq C_{\underline h,\overline h} \mathcal P(\norm{V}{L^\infty}) \norm{\partial_x V}{L^\infty}, \\
    & \norm{\partial_t \Psi}{L^\infty} \leq (\norm{\Psi}{L^{\infty}}+\norm{ \partial_x \Psi }{L^\infty}) \norm{\partial_t V}{L^\infty} \leq C_{\underline h,\overline h} \mathcal P(\norm{V}{L^\infty}) \norm{\partial_t V}{L^\infty}, \\
    & \label{est:nonlinearity-003} \begin{aligned}& \norm{\partial_{x}\partial_t \Psi}{L^2} \leq \norm{\partial_x \Psi}{L^\infty} \norm{\partial_{x}\partial_t V}{L^2} + \norm{\partial_{x}\partial_t\Psi}{L^\infty} \norm{\partial_x V}{L^\infty}\norm{\partial_t V}{L^2} \\
    & \qquad \leq C_{\underline h, \overline h} \mathcal P(\norm{V}{L^\infty}) (\norm{\partial_{x}^2 V}{L^2} + \norm{\partial_x V}{L^\infty}\norm{\partial_t V}{L^2} ),
    \end{aligned}\\
    & \label{est:nonlinearity-004} \begin{aligned}& \norm{\partial_{t}^2 \Psi}{L^2} \leq \norm{\partial_x \Psi}{L^\infty} \norm{\partial_{t}^2 V}{L^2} + \norm{\partial_{x}^2\Psi}{L^\infty} \norm{\partial_t V}{L^\infty}\norm{\partial_t V}{L^2} \\
    & \qquad \leq C_{\underline h, \overline h} \mathcal P(\norm{V}{L^\infty}) (\norm{\partial_{t}^2 V}{L^2} + \norm{\partial_t V}{L^\infty}\norm{\partial_t V}{L^2} ),
    \end{aligned}
\end{align}
where, hereafter, $\mathcal{P}(\cdot) $ is a generic non-decreasing non-negative polynomial of the argument, different from line to line. Therefore, after applying the Sobolev embedding inequality, we arrive at
\begin{equation}
    \label{est:nonlinear-005}
    \sum_{\Psi\in\lbrace A_0, A_1, S_\phi \rbrace}\lbrace \norm{\partial_x \Psi, \partial_t \Psi}{L^\infty} + \norm{ \partial_{x}\partial_t \Psi,\partial_{t}^2 \Psi}{L^2} \rbrace \leq  \mathcal P(\mathcal E(t)^{\frac{1}{2}}).
\end{equation}
Now we conclude with
\begin{equation}
    \label{est:nonlinear-006}
    \begin{gathered}
    \norm{S_\phi,G_{tt},G_{xt},G_{xx}}{L^2}  \leq \sum_{\Psi \in \lbrace A_0, A_1, S_\phi\rbrace} \bigl\lbrace \norm{\partial_x \Psi, \partial_t \Psi}{L^\infty} \norm{\partial_{x}\partial_tV,\partial_{t}^2V,\partial_{x}^2V}{L^2} \\
    +  \norm{\partial_{x}\partial_t \Psi, \partial_{t}^2 \Psi}{L^2} \norm{\partial_{x}V,\partial_{t}V}{L^\infty} \rbrace + \norm{S_\phi, \partial_{x}^2S_\phi, \partial_{x}\partial_t S_\phi, \partial_{t}^2 S_\phi}{L^2} 
    \leq  \mathcal P(\mathcal E(t)^{\frac{1}{2}}).
    \end{gathered}
\end{equation}
Hence, from \eqref{h-002} and \eqref{est:nonlinear-006}, thanks to \eqref{energy-eqvlt}, there exists $T>0$ such that
\begin{equation}
    \label{energy-ineq-total}
    \dfrac{d}{dt}\mathfrak E(t) \leq  \mathcal P(\mathfrak E(t)^{\frac{1}{2}}),
\end{equation}
for all $t\in (0,T)$. This finishes the proof of Proposition \ref{prop:a-priori-V} after integrating \eqref{energy-ineq-total}.


\subsection{Justifying the \textit{a priori} assumption \eqref{non-degenerate}}
\label{Justification}
We now turn to the assumption \eqref{non-degenerate}. Our goal is to show the following proposition:
\begin{proposition}
    \label{prop:a-priori-assumption}
    Suppose that 
    \begin{equation}
    \label{est:energy}
        \sup_{0\leq t \leq T} \mathcal E(t) < C_{\underline h, \overline h} \mathcal E(0),
    \end{equation}
    for some $ T \in(0,\infty) $ and $ C_{\underline h, \overline h} \in (0,\infty) $. There exists $ T'\in (0,T] $ such that 
    \begin{equation}
    \label{est:height}
       0<\frac{3}{4} \underline h \leq \inf_{x,t\in [0,T']} \eta(x,t) \leq \sup_{x,t\in [0,T']} \eta(x,t) \leq \frac{3}{2} \overline h.
    \end{equation}
    In particular, \eqref{est:height} is a stronger estimate than \eqref{non-degenerate}.
    Therefore, after updating $ T $ with $ T' $, both \eqref{est:energy} and \eqref{est:height} hold in the interval $ [0,T] $. 
\end{proposition}


\begin{proof}
    Given $ t \in [0,T] $,
    \begin{equation}
    \label{est:h-001}
        \norm{\partial_t \eta(t)}{L^\infty} \leq \norm{\partial_t \eta(t),\partial_{x}\partial_t \eta(t)}{L^2} \lesssim \mathcal E(t)^{\frac{1}{2}} \lesssim \mathcal E(0)^{\frac{1}{2}},
    \end{equation}
    by the Sobolev embedding inequality. One has that, for all $ t \in [0,T] $,
    \begin{equation}
    \label{est:h-002}
    \begin{gathered}
    \eta(x,t)  = \eta_0(x) + \int_0^t \partial_t \eta(x,\tau) d\tau \leq \eta_0(x) + \int_0^t \norm{\partial_t \eta(\tau)}{L^\infty} \,d\tau\\
    \leq \overline h + t C \mathcal E(0)^{\frac{1}{2}},
    \end{gathered}
    \end{equation}
    for some $ C \in (0,\infty) $. 
    Therefore, by choosing $ t $ small enough, \eqref{est:h-002} implies the last inequality in \eqref{est:height}. The other inequality in \eqref{est:height} follows similarly. 
\end{proof}

We henceforth conclude from Propositions \ref{prop:a-priori-V} and \ref{prop:a-priori-assumption} via boostrap arguments with the following:
\begin{proposition}
    \label{EnergyProposition}
    Consider $ (\eta_0 -1, \omega_0, \phi_0 ) \in H^3(\mathbb{R}) \times H^4(\mathbb{R}) \times H^4(\mathbb{R}) $, i.e., initial data as in \eqref{V-initialdata}. 
    There exists some time $T>0$ such that if $(\eta-1,\omega,\phi)$ solves \eqref{ShallowWaterV}, then
    \begin{equation}
    \label{est:energy-total}
        \sup_{0\leq t\leq T} \sum_{\substack{i+j\leq 2,\\ i,j \geq 0}} \Vert\partial_t^i\partial_x ^j (\eta-1),\partial_t^i\partial_x^j \omega,\partial_t^i\partial_x^{j}\phi\Vert_{L^2}^2 \leq C_{\underline h, \overline h} \mathcal E(0),
    \end{equation}
    for some constant $ C_{\underline h, \overline h} \in (0,\infty) $ depending only on $ \underline h, \overline h $. 
\end{proposition}


\section{Local well-posedness of the 1d generalized system \eqref{ShallowWaterV} via vanishing viscosity}
\label{vv}

\subsection{Vanishing viscosity and the approximating system}

In this section, we will focus on the construction of the local-in-time solution and establish the well-posedness of system \eqref{ShallowWaterV}. Notice that the system \eqref{ShallowWaterV} becomes singular in $ \lbrace h = 0 \rbrace $. To construct a solution to the nonlinear problem, one will need to avoid this possible degeneracy. 
For $ \varepsilon \in (0,1) $, let 
\begin{equation}
\label{def:jpn-bracket}
\langle (\cdot)\rangle_{\varepsilon}:=\sqrt{(\cdot)^2+\varepsilon^2}.
\end{equation}
and 
\begin{equation}\label{vsc-As}\begin{gathered}
A_0^{\varepsilon}:=\begin{pmatrix}
1 & 0 & 0 \\
0 & \langle \eta^\varepsilon\rangle_{\varepsilon}^{-1} & 0 \\
0 & 0 & 0\\
\end{pmatrix}, \quad 
A_1^{\varepsilon}:=
\begin{pmatrix}
    0 & 1 & 0 \\
    1 & 2\omega^\varepsilon\langle \eta^\varepsilon\rangle_{\varepsilon}^{-2} & 0 \\
    0 & 0 & 0 \\
\end{pmatrix},
\\
A_2^{\varepsilon}:=\begin{pmatrix}
        -\varepsilon & 0 & 0 \\
        0 & -\varepsilon & -1 \\
        0 & 1 & -\varepsilon \\
    \end{pmatrix},
    \quad
    S_{\phi}^{\varepsilon}:=
    \begin{pmatrix}
        0 & \\
        (\omega^\varepsilon)^2 \langle \eta^\varepsilon\rangle_{\varepsilon}^{-3}\phi^\varepsilon\\
        0
    \end{pmatrix}.
\end{gathered}
\end{equation}
Then we consider the following regularization of system \eqref{ShallowWaterV}:
\begin{subequations}
\label{ShallowWaterMomentumgeneralized}
\begin{gather}
\label{ShallowWaterMomentumgeneralizedConsice}
A_0^\varepsilon \partial_t V^\varepsilon + A_1^\varepsilon \partial_x V^\varepsilon + A_2^\varepsilon \partial_{x}^2 V^\varepsilon = S_\phi^\varepsilon, \\
%     \begin{pmatrix}
%         1 & 0 & 0 \\
%         0 & \langle h\rangle_{\varepsilon}^{-1} & 0 \\
%         0 & 0 & 1 \\
%     \end{pmatrix}\partial_t V+\begin{pmatrix}
%         0 & 1 & 0 \\
%         1 & 2m \langle h \rangle_{\varepsilon}^{-2} & 0 \\
%         0 & 0 & 0 \\
%     \end{pmatrix}\partial_x V+\begin{pmatrix}
%         -\varepsilon & 0 & 0 \\
%         0 & -\varepsilon & -1 \\
%         0 & 1 & -\varepsilon \\
%     \end{pmatrix}\partial_x^2 V=\begin{pmatrix}
%         0 \\
%         {m^2}\langle h\rangle_{\varepsilon}^{-3}\phi \\
%         0 \\
%     \end{pmatrix}, \\
    V^\varepsilon(0,x)=V_0^\varepsilon = (\eta_0^{\varepsilon}-1, \omega_0^{\varepsilon}, \phi_0^{\varepsilon}), 
\end{gather}
\end{subequations}
where the initial data satisfies 
\begin{equation}
    \label{Vsc_initial}
    \begin{aligned}
    \eta_0^\varepsilon - 1 & \xrightarrow{\varepsilon\to 0} \eta_0-1  \qquad  & \text{in} ~ H^3(\mathbb R), \\
    (\omega_0^{\varepsilon}, ~ \phi_0^{\varepsilon}) & \xrightarrow{\varepsilon\to 0} (\omega_0, ~ \phi_0) & \text{in} ~ H^4(\mathbb R),
    \end{aligned}
\end{equation}
and 
\begin{align}
    \label{Vsc_initial-2}
    \varepsilon (\eta_0^\varepsilon-1) & \xrightarrow{\varepsilon\to 0} 0  \qquad   \text{in} ~ H^4(\mathbb R). 
\end{align}
The energy associated to \eqref{ShallowWaterMomentumgeneralized} is given by
\begin{equation}
    \label{EpsilonEnergy}
    \mathfrak E^{\varepsilon}(t):=\sum_{\substack{i+j\leq 2,\\ i,j \geq 0}} \Vert\partial_t^i\partial_x^j (\eta^{\varepsilon}-1),\frac{\partial_t^i\partial_x^j \omega^{\varepsilon}}{{\langle \eta^{\varepsilon}\rangle}^{\frac{1}{2}}_{\varepsilon}},\partial_t^i\partial_x^{j}\phi^{\varepsilon}\Vert_{L^2}^2.
\end{equation}
For any fixed $ \varepsilon $, system \eqref{ShallowWaterMomentumgeneralized} is a non-degenerate uniformly parabolic system, for which there exists a local-in-time solution through construction in, for instance, \cite{Taylor}. Note that the existence time of the said solution may depend on $ \varepsilon $. However, one can immediately check that the {\it a priori} estimates in Section \ref{1DEnergyEstimates} still hold for $ V^\varepsilon $, uniformly in $ \varepsilon $. Therefore, one can conclude the following proposition:
\begin{proposition}
    \label{WellPosedViscousProblem}
    Given initial data $ V^\varepsilon_0 = (\eta_0^{\varepsilon}-1,\omega_0^{\varepsilon}, \phi_0^{\varepsilon})\in H^4(\mathbb{R})\times H^4(\mathbb{R})\times H^4(\mathbb{R})$, 
    there exists $ \mathfrak{T}\in (0,\infty) $, independent of $ \varepsilon $, such that there exists a unique solution $ V^\varepsilon = (\eta^\varepsilon-1,\omega^\varepsilon,\phi^\varepsilon) \in  L^\infty(0,\mathfrak{T};H^2(\mathbb{R})) $ and $ \partial_t V^\varepsilon \in L^\infty(0,\mathfrak{T};H^1(\mathbb R)) $, satisfying
    \begin{equation}
        \label{energy-viscous}
        \sup_{0\leq t \leq \mathfrak{T}} \sum_{\substack{i + j \leq 2, \\
        i \geq 0, j \geq 0}} \norm{\partial_t^i\partial_x^j V^\varepsilon(t)}{L^2} \leq C (\varepsilon\norm{ \eta^\varepsilon_0-1}{H^4}+\norm{\eta_0^\varepsilon-1}{H^3} + \norm{\omega^\varepsilon_0,\phi^\varepsilon_0}{H^{4}} ),
    \end{equation}
    for some constant $ C \in (0,\infty) $, depending only on $ \underline h, \overline h $ but independent of $ \varepsilon \in (0,1) $. 
\end{proposition}

\begin{proof}
From Proposition \ref{EnergyProposition}, we only need to verify the initial data. In particular, thanks to similar expansion of the initial data as in \eqref{initial_data_time}, the only item we need to verify is $\varepsilon \norm{ \eta_0^{\varepsilon}-1}{H^4}$. However, notice that for the viscous system \eqref{ShallowWaterMomentumgeneralizedConsice}, we have
\[
\partial_t \eta^{\varepsilon}(t=0)=\varepsilon \partial_x^2 \eta^{\varepsilon}_0-\partial_x\omega^{\varepsilon}_0,
\]
and
\begin{align*}
    \partial_t^2 \eta^{\varepsilon}(t=0) &= -2\frac{\omega_0^{\varepsilon} }{(\eta_0^{\varepsilon})^2} \phi_0^{\varepsilon}\partial_x \omega_0^{\varepsilon}-\frac{(\omega_0^{\varepsilon})^2}{(\eta_0^{\varepsilon})^2 }\partial_x \phi_0^{\varepsilon}+\frac{2(\omega_0^{\varepsilon})^2}{(\eta_0^{\varepsilon})^3}\phi^{\varepsilon}_0\partial_x\eta^{\varepsilon}_0-\partial_x\eta_0^{\varepsilon} \partial_x^3\phi_0^{\varepsilon}-\eta_0^{\varepsilon}\partial_x^3 \phi_0^\varepsilon+(\partial_x\eta_0^{\varepsilon})^2  \\
    & +\eta_0^{\varepsilon}\partial_x^2 \eta^{\varepsilon}_0+\frac{2(\partial_x\omega_0^{\varepsilon})^2}{\eta_0^{\varepsilon}}+\frac{2\omega_0^{\varepsilon}}{\eta_0^{\varepsilon}}\partial_x^2 \omega_0^{\varepsilon}-\frac{2\omega_0^{\varepsilon}}{(\eta_0^{\varepsilon})^2 } \partial_x \eta^{\varepsilon}_0\partial_x \omega^{\varepsilon}_0-\varepsilon\partial_x^4\eta_0^{\varepsilon}, \\
\end{align*}
which, comparing to \eqref{initial_data_time}, implies that we need one more degree of space regularity for $\eta^{\varepsilon}_0-1$ than that of $\eta_0-1.$ %By the same reasoning, the term $\varepsilon\partial_x^4 \eta_0^{\varepsilon}$ appears in the energy $\mathfrak{E}^{\varepsilon}(t)$ associated to \eqref{ShallowWaterMomentumgeneralizedConsice}. 
Analogously to \eqref{est:energy}, one can show
\[
\mathfrak E^{\varepsilon}(t)\leq 2\mathfrak E^{\varepsilon}(0).
\]
This in turn implies
\[
\mathfrak E^{\varepsilon}(t)\lesssim \norm{\partial_t \eta_0^{\varepsilon}}{H^1}^2+\norm{\partial_t^2 \eta_0^{\varepsilon}}{L^2}^2+\norm{\omega_0^{\varepsilon}, \phi_0^{\varepsilon}}{H^4}^2\lesssim \varepsilon^2 \norm{\eta_0^{\varepsilon}-1}{H^4}^2+\norm{\eta_0^{\varepsilon}-1}{H^3}^2+\norm{\omega^{\varepsilon}_0,\phi_0^{\varepsilon}}{H^4}^2<\infty.
\]
% thanks to the modes of convergence \eqref{Vsc_initial} and \eqref{Vsc_initial-2}.
Here $ \mathfrak E^\varepsilon $ is the energy for the approximating system \eqref{ShallowWaterMomentumgeneralized} in analogy to \eqref{intrmd-energy}. This finishes the proof. 
\end{proof}

\subsection{Existence of solutions}
\label{Ex-U}
With the uniform estimate \eqref{energy-viscous}, together with \eqref{Vsc_initial} and \eqref{Vsc_initial-2}, one can pass the limit $ \varepsilon \rightarrow 0 $ and, up to a subsequence, obtain
\begin{align}
    (\eta^\varepsilon - 1, ~ \omega^\varepsilon,~ \phi^\varepsilon) & \overset{*}{\rightharpoonup}  (\eta -1, ~ \omega, ~ \phi) & \text{weakly-$*$ in} ~ L^\infty([0,\mathfrak{T}];H^2(\mathbb R)),
\end{align}
for the same $\mathfrak{T}$ given in Proposition \ref{WellPosedViscousProblem}. By the Aubin-Lions compactness lemma, 
\begin{align}
\label{1-d-cnv:strong}
    (\eta^\varepsilon - 1, ~ \omega^\varepsilon, \phi^\varepsilon) & \rightarrow  (\eta -1, ~ \omega, ~ \phi) & \text{strongly in} ~ C([0,\mathfrak{T}];H^1_{loc}(\mathbb R)),
\end{align}
for some $ V = (\eta-1,\omega,\phi) \in L^\infty([0,\mathfrak{T}];H^2(\mathbb R)) $, satisfying 
\begin{equation}
        \label{energy-viscous-2}
        \sup_{0\leq t \leq T} \sum_{\substack{i + j \leq 2, \\
        i \geq 0, j \geq 0}} \norm{\partial_t^i\partial_x^j V(t)}{L^2} \leq C (\norm{\eta_0-1}{H^3} + \norm{\omega_0,\phi_0}{H^{4}} ),
    \end{equation}
thanks to \eqref{Vsc_initial} and \eqref{Vsc_initial-2}. Here $ H_{loc}^1(\mathbb R) $ representing the local $ H^1 $ space. That is, the strong convergence \eqref{1-d-cnv:strong} holds inside arbitrary compact subset of $ \mathbb R $. 

It is straightforward to verify that $ V $ solves \eqref{ShallowWaterV} with the above weak and strong compactness. 

\iffalse

By Proposition \ref{EnergyProposition}, there exists a uniform time $\mathfrak{T}\in (0, \mathcal{T}(\varepsilon)]$ independent of $\varepsilon$ such that 
\begin{equation}
\label{UniformTime}
\sup_{0\leq t\leq \mathfrak{T}}E^{\varepsilon}(t)\leq \sup_{0\leq t\leq \mathfrak{T}} E(t)\leq 1.
\end{equation}
Hence, by Proposition \ref{WellPosedViscousProblem} and the convergence modes \eqref{Vsc_initial}, \eqref{Vsc_initial-2}, for the solution $U_*^{\varepsilon}$ to the viscosity problem \eqref{ShallowWaterMomentumgeneralized}, there exists a subsequence $\varepsilon_n\to 0$ such that 
\[
U_*^{\varepsilon_n}\rightarrow U_* \quad \text{uniformly in} \quad L^2([0,\mathfrak{T}];H^2(\mathbb{R}))\cap L^{\infty}([0,\mathfrak{T}];L^2(\mathbb{R})),
\]
for the above choice of $\mathfrak{T}$. In particular, by the argument in Section \ref{Justification}, the height $h$ is strictly bounded away from zero, and so the convergence of $\langle h^{\varepsilon}\rangle_{\varepsilon}^{-a}$ is not an issue. The limit $U_*\in L^2([0,\mathfrak{T}];H^2(\mathbb{R}))\cap L^{\infty}([0,\mathfrak{T}];L^2(\mathbb{R}))$ solves \eqref{ShallowWaterMomentum} with initial data $V_*(0,x)=(h_0,m_0,\phi_0)$.
\fi


\subsection{Uniqueness and continuous dependence on initial data}
\label{subsec:well-posedness-generalized}

% Denote $\mathcal{V}$ to be the above solution to \eqref{ShallowWaterV} obtained from vanishing viscosity. 
Let $\mathcal V, \tilde{\mathcal{V}} \in L^{\infty}([0,\mathfrak{T}];H^2(\mathbb{R}))$ be two solutions to \eqref{ShallowWaterV} with initial data $\mathcal V(0,x) = \mathcal V_0:= (\eta_0-1,\omega_0,\phi_0), \tilde{\mathcal{V}}(0,x) = \tilde{\mathcal V}_0 :=(\tilde{\eta}_0-1,\tilde{\omega}_0,\tilde{\phi}_0)\in H^3(\mathbb{R}) \times H^4(\mathbb{R}) \times H^4(\mathbb{R})$, respectively. Then $\delta\mathcal{V}:=\mathcal{V}-\tilde{\mathcal{V}}=(\delta\eta,\delta\omega,\delta\phi)^\top$ solves
\[
\begin{cases}
A_0(\tilde{\mathcal{V}})\partial_t\delta\mathcal{V}+A_1(\tilde{\mathcal{V}}) \partial_x \delta\mathcal{V}+A_2 \partial_x^2 \delta \mathcal{V}=\begin{pmatrix}
    0 \\ 
    N(\mathcal{V},\tilde{\mathcal{V}}) \\
    0 \\
\end{pmatrix},\\
    \delta\mathcal{V}(0,x)=\delta\mathcal{V}_0=(\eta_0-\tilde{\eta}_0,\omega_0-\tilde{\omega}_0,\phi_0-\tilde{\phi}_0),
\end{cases}
\]
with
\[
N(\mathcal{V},\tilde{\mathcal{V}}):=\frac{\omega^2}{\eta^3} \phi-\frac{\tilde{\omega}^2}{\tilde{\eta}^3} \tilde{\phi}-\underbrace{[A_0(\mathcal{V})-A_0(\tilde{\mathcal V})]}_{\in L^{2}}\underbrace{\partial_t \mathcal V}_{\in L^\infty}-\underbrace{[A_1(\mathcal{V})-A_1(\tilde{\mathcal V})]}_{\in L^{2}}\underbrace{\partial_x \mathcal V}_{\in L^\infty}.
\]
Multiplying by $\delta\mathcal{V}^\top$ and integrating in space, the $L^2$ energy estimate for $\delta \mathcal V $ reads
\begin{equation*}
    \int \delta\mathcal{V}^\top A_0(\tilde{V})\partial_t \delta\mathcal{V}+\delta\mathcal{V}^\top A_1(\tilde{V}) \partial_x \delta\mathcal{V}+\delta\mathcal{V}^\top A_2 \partial_x^2 \delta\mathcal
    V dx=\int(\omega-\tilde{\omega}) N(\mathcal{V},\tilde{\mathcal{V}}) dx.
\end{equation*}
Denote $\tilde{\mathcal{E}}(t)$ the sum of the energy functionals for $\mathcal{V}$ and $\tilde{\mathcal{V}}$. A direct estimate yields
\begin{align*}
& \frac{1}{2}\frac{d}{dt}\norm{\delta\mathcal{V}}{L^2}^2\lesssim \tilde{\mathcal{P}}(\tilde{\mathcal{E}}(t)^{\frac{1}{2}}) \norm{\delta \mathcal{V}}{L^2}^2,
\end{align*}
for some positive polynomial $\tilde{\mathcal{P}}$. Now Grönwall's inequality implies 
\begin{equation}
\label{ContDepend}
\norm{\delta\mathcal{V}}{L^2}^2\lesssim \exp \left(\int_0^{\mathfrak{T}} \tilde{\mathcal{P}}(\tilde{\mathcal{E}}(t)^{\frac{1}{2}}) dt\right) \norm{\delta\mathcal{V}_0}{L^2}^2.
\end{equation}
In particular, if $\mathcal{V}_0=\tilde{\mathcal{V}}_0$, then $\mathcal{V}\equiv \tilde{\mathcal{V}}.$ This finishes the proof of well-posedness of the generalized system \eqref{ShallowWaterV}.

\section{Well-posedness of the 1d shallow water system \eqref{ShallowWaterMomentumConsice}}
\label{Sect:WP}


We now are ready to establish well-posedness system \eqref{ShallowWaterMomentumConsice} for initial data with regularity $(h_0-1,m_0)\in H^3(\mathbb R)\times H^2 (\mathbb R)$. To do so, we need elliptic estimates and energy estimates for the original shallow water system \eqref{ShallowWaterMomentumConsice}. We emphasize that the elliptic estimates are not available for the generalized system \eqref{ShallowWaterV}; see Remark \ref{rm:non-elliptic-est-general-sys}, below. 
\subsection{Elliptic estimates} We define the evolutionary energy 
\begin{equation}
    \label{EnergyEvolution}
    \mathcal F(t) := \sum_{0\leq i\leq 1} \norm{\partial_t^i (h-1),\frac{\partial_t^i m}{\sqrt h}, \partial_t^i \partial_x h}{L^2}^2,
\end{equation}
and the elliptic energy 
\begin{equation}
\label{EnergyElliptic}
    \mathfrak F(t):=\sum_{0\leq j\leq 2} \Vert \partial_x ^j (h-1),\frac{\partial_x^j m}{\sqrt{h}},\partial_x^{j}\partial_x h\Vert_{L^2}^2.
\end{equation}
Similar to \eqref{non-degenerate}, we make the {\it a priori} assumption for $ h $:
\begin{equation}
    \label{non-degenerate-h}
  0<\frac{1}{2}\underline h \leq \inf_{x,t} h(x,t) \leq \sup_{x,t} h(x,t) \leq 2 \overline h.
\end{equation}
We make the following observation:
\begin{lemma}
    \label{EllipticLemma} 
    Under the {\it a priori} assumption \ref{non-degenerate-h},
   there exists a positive polynomial $\mathcal{G}_1$ %and $\mathcal{G}_2$ 
   such that
   \begin{equation}
   \label{EllipticEstimate1}
       \mathfrak F(t)\leq \mathcal{G}_1(\mathcal{F}(t)).
       % , \\
       % \label{EllipticEstimate2}
       % &\frac{d}{dt} \mathcal{F}(t)\leq \mathcal{G}_2(\mathfrak F(t)).
   \end{equation}

\begin{proof}
For the original system \eqref{ShallowWaterMomentumConsice}, we have
        \begin{gather}
        \label{EllipticHeight}
            \norm{\partial_x m}{L^2}= \norm{\partial_t h}{L^2}, \\
            \label{EllipticHeightD}
            \norm{\partial_x^2 m}{L^2}=\norm{\partial_t \partial_x h}{L^2}.
        \end{gather}
    At the same time, applying integration by parts, one has that
    \begin{equation}
            \label{EllipticMomentum}
    \begin{gathered}
            \norm{\partial_x^3 h}{L^2}^2+\norm{\partial_x h}{L^2}^2 + 2 \norm{\partial_x^2 h}{L^2}^2 = \norm{-\partial_{x}^3h + \partial_x h}{L^2}^2\\
            = \norm{\frac{1}{h}\partial_t m + \frac{2 m \partial_x m}{h^2} - \frac{m^2 \partial_x h}{h^3}}{L^2}^2 \lesssim \norm{\partial_t m}{L^2}^2+\norm{m}{H^1}^4+\norm{m}{H^1}^4 \norm{\partial_x h}{L^2}^2,
        \end{gathered}
    \end{equation}
    where we have applied \eqref{non-degenerate-h}.
This finishes the proof of \eqref{EllipticEstimate1}.

% The identities \eqref{EllipticHeight}--\eqref{EllipticMomentum} follow from \eqref{ShallowWaterMomentum}, the Cauchy--Schwarz inequality, the regularity prescribed in Section \ref{Ex-U}, and the assumption \eqref{non-degenerate}. Subsequently, to have $\partial_t U\in L^2_x(\mathbb{R})$, we need
% \begin{equation}
% \label{EllipticCorollary}
% h-1\in H^3_x(\mathbb{R}) \quad \text{and} \quad m\in H_x^2(\mathbb{R}).
% \end{equation}
% Moreover, by \eqref{EllipticHeight}--\eqref{EllipticHeightD}, we have
% \[
% \norm{\partial_x h}{L^2}+\norm{\partial_x^3 h}{L^2}+\norm{\partial_x m}{L^2}+\norm{\partial_x^2 m}{L^2}\leq \mathcal{F}(t)+\mathcal{F}(t)^2+\mathcal{F}(t)^3,
% \]
% and by choosing $\mathcal{G}_1$ to be a positive cubic polynomial, we obtain \eqref{EllipticEstimate1} thanks to the Sobolev interpolation inequality
% \[
% \norm{\partial_x^2 h}{L^2}\leq \norm{\partial_x h}{L^2}^{\frac{1}{2}} \cdot \norm{\partial_x^3 h}{L^2}^{\frac{1}{2}}.
% \]
% Lastly, \eqref{EllipticEstimate2} follows from Lemma \ref{EnergyOriginal} and Proposition \ref{apriori-Original}.
\end{proof}
\end{lemma}
\begin{remark}
\label{rm:non-elliptic-est-general-sys}
    Lemma \ref{EllipticLemma} does not hold for the generalized system \eqref{ShallowWaterV}. Indeed, in the $(\eta-1,\omega,\phi)$-variable, the corresponding estimate \eqref{EllipticMomentum} becomes 
    \[
    \norm{-\partial_x^2 \phi+\partial_x h}{L^2}^2 \lesssim \norm{\partial_t \omega}{L^2}^2+\norm{\omega}{H^1}^4+\norm{\omega}{H^1}^2\norm{\phi}{L^2}.
    \]
    In particular, since $ \phi \neq \partial_x h $ in general, one cannot obtain the higher order regularity of $ h $ from the left hand side of the above estimate.

    By making use of the elliptic estimate \eqref{EllipticEstimate1} for the shallow water system \eqref{ShallowWaterMomentumConsice}, we will be able to obtain sharper regularity for the solutions.
    
% As a result, we cannot determine the regularity of $\partial_x h$ and $\partial_x^2 \phi$ individually. This is no longer an issue in \eqref{EllipticMomentum}, for which $\eta=h$, $\phi=\partial_x h$ and $h-1\in H^3_x(\mathbb{R})$ trivially implies $h-1\in H^1_x(\mathbb{R}).$
\end{remark}
\subsection{Energy estimates} For the shallow water system \eqref{ShallowWaterMomentumConsice}, we have the following analog to Lemma \ref{LemmaPhi}:
\begin{lemma}
\label{EnergyOriginal}
    Let $\Phi:\mathbb R \times\mathbb R\to\mathbb{R}^3$ be the solution to the linear equation
    \[
    A_0(U)\partial_t \Phi +A_1(U)\partial_x \Phi+A_2 \partial_x^2 \Phi=G.
    \]
    Then for a.e. $t$, we have
    \begin{equation}
        \frac{d}{dt} (A_0(U)\Phi,\Phi)\leq C_{\underline{h},\overline{h}} ( \mathcal{F}(t) + \mathfrak F(t) ) \norm{\Phi}{L^2}^2+2\norm{G}{L^2} \norm{\Phi}{L^2}.
    \end{equation}
% \begin{proof}
% The proof of Lemma \ref{EnergyOriginal} is near identical to that of Lemma \ref{LemmaPhi}, except the last step \eqref{phi:003} is replaced by
% \[
% \norm{\partial_t A_0,\partial_x A_1}{L^{\infty}}\leq C_{\underline h} \sum_{j=0}^2 \norm{\partial_x^j m,\partial_x^j (h-1)}{L^2}^2\leq C_{\underline{h},\overline{h}} \mathcal{F}(t),
% \] via the elliptic estimate \eqref{EllipticEstimate1}.
% \end{proof}
\end{lemma}
The proof of Lemma \ref{EnergyOriginal} follows with the same arguments as in Lemma \ref{LemmaPhi}.

\smallskip 

One can similarly prove the analog of Proposition \ref{prop:a-priori-V} for $\mathcal{F}.$ That is:
\begin{proposition}
    \label{apriori-Original}
    % Provided that \eqref{non-degenerate-h} holds, 
    The solution to \eqref{ShallowWaterMomentumConsice} satisfies the uniform bound
    \[
    \sup_{0\leq t\leq {T}} \mathcal{F}(t)\leq 2 \mathcal{F}(0),
    \]
    for some $T\in (0,\infty)$ depending only on the initial data and $\underline h$, $\overline{ h}.$
\end{proposition}

\begin{proof}
    Proposition \ref{apriori-Original} follows from an argument similar to that in Section \ref{1DEnergyEstimates}. Here we only sketch the main steps. We only consider, for the shallow water system \eqref{ShallowWaterMomentumConsice}, 
    \begin{equation}
    \label{sys:dd-Original}
    A_0(U)\partial_t U_t+A_1(U)\partial_x U_t+A_2\partial_{x}^2 U_t=\mathfrak G_t,
    \end{equation}
    where 
    \begin{align*}
        & \mathfrak G_t:=\partial_t S(U)-\partial_t A_0(U) \partial_t U-\partial_t A_1(U)\partial_x U.
    \end{align*}
    Applying Lemma \ref{EnergyOriginal} to systems \eqref{ShallowWaterMomentumConsice} and \eqref{sys:dd-Original}, we obtain
\begin{equation}
    \label{h2:001-original}
    \frac{d}{dt} \mathcal F(t) \leq C_{\overline {h},\underline{h}} (\mathcal{F}(t) + \mathfrak F(t)) \norm{U,U_{t}}{L^2}^2 +2\norm{{S(U)},\mathfrak{G}_{t}}{L^2}\norm{U,U_{t}}{L^2} \leq \mathcal G_2(\mathcal F(t)),
\end{equation}
for some positive polynomial $ \mathcal G_2 $, thanks to Lemma \ref{EllipticLemma} and embedding inequalities. Integrating \eqref{h2:001-original} finishes the energy estimate. Moreover, following the argument as in Section \ref{Justification}, one can establish the justification of \eqref{non-degenerate-h}. This finishes the proof. 
\end{proof}
%  Here we point out the main differences: For the original system \eqref{ShallowWaterMomentumConsice}, we consider the systems
% \begin{equation}
%     \label{sys:dd-Original}
%     A_0(U)\partial_t U_a+A_1(U)\partial_x U_a+A_2\partial_{x}^2 U_a=\mathfrak G_a \quad \text{for}\quad a\in \{t,xt,xx\}, 
% \end{equation}
% where 
% \begin{align*}
%     & \mathfrak G_t:=\partial_t S(U)-\partial_t A_0(U) \partial_t U-\partial_t A_1(U)\partial_x U, 
% \end{align*}
% % and $\mathfrak G_{xt}$, $\mathfrak G_{xx}$ are analogous to \eqref{id:xt_G} and \eqref{id:xx_G}. Applying Lemma \ref{apriori-Original} to \eqref{sys:dd-Original}, we obtain
% \begin{equation}
%     \label{h2:001-original}
%     \frac{d}{dt}\sum_{a\in \{t,xt,xx\}} (A_0(U)U_a,U_a) \leq C_{\overline {h},\underline{h}} \mathcal{F}(t) \norm{U_t,U_{xt},U_{xx}}{L^2}^2 +2\norm{\mathfrak{G}_t,\mathfrak{G}_{xt},\mathfrak{G}_{xx}}{L^2}\norm{U_t,U_{xt},U_{xx}}{L^2}.
% \end{equation}
% Similarly to the estimates \eqref{h-1-eq}--\eqref{est:nonlinear-005}, we arrive at 
% \begin{align*}
%  & \norm{S(U),\mathfrak G_t,\mathfrak G_{xt}, \mathfrak G_{xx}}{L^2} \\ 
%  & \leq \sum_{\{\Psi\in A_0(U),A_1(U),S\}} \{\norm{\partial_x \Psi,\partial_{t}\Psi}{L^{\infty}} \norm{\partial_t U,\partial_{x}\partial_t U,\partial_{x}^2 U}{L^2}+\norm{\partial_t \Psi, \partial_{x}\partial_t \Psi}{L^2} \norm{\partial_x U,\partial_T U}{L^2} \} \\
%  &+\norm{S(U),\partial_x^2 \phi,\partial_x\partial_t S(U),\partial_t S(U)}{L^2}\leq C_{\underline{h},\overline{h}} \mathfrak{P}(\mathcal{F}(t)^{\frac{1}{2}}), 
% \end{align*}
% for some positive polynomial $\mathfrak P$. Now Proposition \ref{apriori-Original} follows from Grönwall's inequality.
\subsection{Proof of Theorem \ref{thm:1-d}}
By taking $\eta=h$, $\omega=m$, and $\phi=\partial_x h$ in system \eqref{ShallowWaterV}, for $(h_0-1,m_0)\in H^5(\mathbb{R})\times H^4(\mathbb{R})$, Proposition \ref{EnergyProposition} implies that there exists a unique solution $ (h-1, m)\in H^3(\mathbb{R})\times H^2(\mathbb{R}) $ to the shallow water system \eqref{ShallowWaterMomentumConsice}.

In general, for $ (h_0-1, m_0) \in H^3(\mathbb R) \times H^2(\mathbb R) $, one chooses a sequence of initial data $ (h_{k,0} - 1, m_{k,0}) \in H^5(\mathbb R) \times H^4(\mathbb R)$, $ k \in \mathbb N $, such that
\begin{equation}
\label{approximation_initial_k}
\begin{aligned}
    h_{k,0} - 1 & \rightarrow h_0 - 1  & \text{in} \ H^3(\mathbb R),\\
    m_{k,0} & \rightarrow m_0 & \text{in} \ H^2(\mathbb R), 
\end{aligned}
\end{equation}
and the initial data
\begin{equation}
    \label{V_k}
    V_{k,0} := \begin{pmatrix}
        h_{k,0} - 1 \\ m_{k,0} \\ \partial_x h_{k,0}
    \end{pmatrix} \in H^3(\mathbb{R}) \times H^4(\mathbb{R}) \times H^4(\mathbb{R}). 
\end{equation}
% satisfies \eqref{V-initialdata}. 
The well-posedness theory of the generalized system \eqref{ShallowWaterV} implies that for each fixed $k=1,2,\ldots$, there exists $ T_k \in (0,\infty) $ such that we have a unique solution $ V_k :=( h_k -1, m_k, \phi_k )^\top $. % satisfying \eqref{energy-viscous-2}. 
In particular, $ (h_k-1, m_k) $ is a solution to the shallow water system \eqref{ShallowWaterMomentumConsice}. 

Applying Proposition \ref{apriori-Original} and Lemma \ref{EllipticLemma}, one obtains the following estimate:
\begin{equation}
\label{EllipticUniform}
\begin{gathered}
\sup_{0\leq t \leq \mathfrak{T}'} \bigl\lbrace 
\norm{h_k(t)-1}{H^3}^2+\norm{m_k(t)}{H^2}^2 + \norm{h_k(t)-1, \partial_t h_k(t)}{H^1}^2 + \norm{m_k(t), \partial_t m(t)}{L^2}^2 \bigr\rbrace \\
\lesssim \norm{\partial_t m_{k,0},m_{k,0}}{L^2}^2+\norm{\partial_t h_{k,0}, h_{k,0}- 1}{H^1}^2 \lesssim  \norm{h_{k,0}-1}{H^3}^2 + \norm{m_{k,0}}{H^2}^2,
\end{gathered}
\end{equation}
for some sufficiently short time $\mathfrak{T}'>0$ independent of $k$. The last inequality follows from directly using the equations \eqref{ShallowWaterMomentumConsice} to obtain the initial data for the time derivatives, similar to \eqref{initial_data_time}.
% Unlike in \eqref{vsc-As} with artificial viscosity, there are no second time derivatives involved in the elliptic estimates \eqref{EllipticHeight}--\eqref{EllipticHeightD}, nor in the energy \eqref{EnergyOriginal}. 
Therefore, by the convergence \eqref{approximation_initial_k} and the bound \eqref{EllipticUniform}, passing the limit $ k \rightarrow \infty $, 
one obtains %only requires $(h_0-1,m_0)\in H^3(\mathbb{R}) \times H^2(\mathbb{R}) $ for the original system \eqref{ShallowWaterMomentum} to have 
the unique solution $(h-1,m)\in L^{\infty}([0,\mathfrak{T}'];H^3(\mathbb R))\times L^{\infty}([0,\mathfrak{T}'];H^2(\mathbb R))$ to the shallow water system \eqref{ShallowWaterMomentumConsice}. The uniqueness and continuous dependence on the initial data follow with arguments similar to those in Section \ref{subsec:well-posedness-generalized}. This concludes the proof of Theorem \ref{thm:1-d} after recalling $u=\frac{m}{h}$.

\section{{\it A priori} estimate for two dimensional flow}
\label{sec:2d} 

\subsection{Symmetric reformulation}
\label{subsec:2d-reformulation}

We will turn our focus to the case of multi-dimensions in system \eqref{ShallowWaterEquations}. In this section, we consider the case when $ n = 2 $, i.e., the two dimensional flow. To be more precise, we will only establish the {\it a priori} estimate using the symmetric structure. The local well-posedness follows via similar arguments as in Sections \ref{vv}--\ref{Sect:WP}.

For $ \mathbb R^2 = \lbrace (x_1, x_2) \rbrace $, we write, for $j =1,2 $,
\begin{equation}
    \label{def:2d-drvt}
    \partial_j := \partial_{x_j}.
\end{equation}
We write, for any vector field $ v = (v_1, v_2)^\top $, 
\[
v^{\perp}:=\begin{pmatrix}
    -v_2 \\
    v_1
\end{pmatrix} = J v, \quad\nabla^{\perp}:= \begin{pmatrix}
    -\partial_2 \\
    \partial_1 \\
\end{pmatrix} = J \nabla ,
\]
where
\begin{equation}
\label{def:anti-matrix}
    J:= \begin{pmatrix}
        0 & -1 \\ 1 & 0
    \end{pmatrix},
\end{equation}
and we consider the 2d flow with the non-degenerate condition
\begin{equation}
    \label{non-degenerate-nD}
    0<\frac{1}{2}\underline h \leq \inf_{x,t} h(x,t) \leq \sup_{x,t} h(x,t) \leq 2 \overline h < \infty. 
\end{equation}
% For the second part of our paper, we prove Theorem \ref{MainTheoremLWP}. We first demonstrate this in dimensions $2$ and $3$, and then generalize to higher dimensions in which $\curl $ is not in general well-defined. For the rest of this paper, we denote $\partial\in \{\partial_1,\ldots,\partial_n\}$ to be an unspecified spatial derivative, and we write
% \[
% m^{\perp}:=\begin{pmatrix}
%     -m_y \\
%     m_x
% \end{pmatrix}, \quad\nabla^{\perp}:= \begin{pmatrix}
%     -\partial_y \\
%     \partial_x \\
% \end{pmatrix}.
% \]
% Similarly to the one-dimensional case, one has to first consider the higher-dimensional case in the $V$-variable \eqref{Def:V}, then go back to the original formulation \eqref{ShallowWaterMomentumConsice} as demonstrated in Section \ref{Sect:WP}. For simplicity, we only show how to close the energy estimates, as the other steps of the proof are similar to the 1D case. In particular, in dimension $n$, we consider the energy functional given by
% \begin{equation}
% \label{Energy_higher_d}
%     \mathcal{E}_n(t):=\sum_{\substack{i+\vert \alpha \vert\leq \lceil \frac{n}{2}\rceil+1,\\ i,\vert\alpha \vert\geq 0}} \Vert\partial_t^i\partial ^{\alpha} (\eta-1),\frac{1}{\sqrt{\eta}}\partial_t^i\partial^{\alpha} \omega,\partial_t^i\partial^{\alpha}\phi\Vert_{L^2}^2
% \end{equation}
% in the $(\eta-1,\omega,\phi)-$variables. Here $\partial^{\alpha}$ represents the usual multi-index notation. We again assume the \emph{a priori} bounds
% \begin{equation}
%     \label{non-degenerate-nD}
%     \frac{1}{2}\underline h \leq \inf_{x,t} \eta(x,t) \leq \sup_{x,t} \eta(x,t) \leq 2 \overline h, 
% \end{equation}
% for all $0\leq t\leq T$, with $T$ being the maximum existence time. One can justify \eqref{non-degenerate-nD} in the same manner as \eqref{non-degenerate}.
% \subsection{The 2D case}
% In dimension 2, we observe that the change of variable $m=hu$ does not directly yield a symmetric system from \eqref{ShallowWaterEquations}. Indeed, the 2D analog of \eqref{ShallowWaterMomentum} reads

% , i.e., 
% \begin{subequations}
% \label{ShallowWaterMomentum2DDivergence}
% \begin{gather}
%     \partial_t h+\Div m=0, \\
%     \label{Momentum2D}
%     \frac{1}{h}\partial_t m+\frac{1}{h}\Div \left(\frac{m\otimes m}{h}\right) +\nabla h-\nabla\Delta h=0, \\
%     \partial_t \nabla h +\nabla \Div m=0.
% \end{gather}
% \end{subequations}


Similarly as before, consider the momentum variable $ m = (m_1,m_2)^\top:= hu = h(u_1,u_2)^\top $. Recall that the shallow water system \eqref{ShallowWaterEquations} in 2d can be written as in \eqref{ShallowWaterMomentum}. We point out, that unlike the 1d case as in \eqref{ShallowWaterMomentumConsice}, the 2d flow \eqref{ShallowWaterMomentum} is not a symmetric system. Indeed, by writing $ U:= (h -1, m, \nabla h)^\top = (h - 1, m_1, m_2, \partial_1 h, \partial_2 h) \in \mathbb R^5 $, one has
\begin{equation}
\label{ShallowWaterMomentum2D}
\diag\left(1,\frac{1}{h},\frac{1}{h},1,1\right) \partial_t U+M_1 \partial_1 U +M_2\partial_2 U +M_3\partial_1^2 U +M_4\partial_2^2 U +M_5\partial_{1}\partial_2 U=S,
\end{equation}
where
\begin{equation}
    \label{def:M-2d}
\begin{aligned}
M_1 & :=\begin{pmatrix}
    0 & 1 & 0 & 0 & 0 \\
    1 & \frac{2m_1}{h^2} & 0 & 0 & 0 \\
    0 & {\frac{m_2}{h^2}} & \frac{m_1}{h^2} & 0 & 0 \\
    0 & 0 & 0 & 0 & 0 \\
    0 & 0 & 0 & 0 & 0 \\
\end{pmatrix}, & M_2 &:=
\begin{pmatrix}
    0 & 0 & 1 & 0 & 0 \\
    0 & \frac{m_2}{h^2} & {\frac{m_1}{h^2}} & 0 & 0 \\
    1 & 0 & \frac{2m_2}{h^2} & 0 & 0 \\
    0 & 0 & 0 & 0 & 0 \\
    0 & 0 & 0 & 0 & 0 \\
\end{pmatrix}, \\
M_3& :=\begin{pmatrix}
    0 & 0 & 0 & 0 & 0 \\
    0 & 0 & 0 & -1 & 0 \\
    0 & 0 & 0 & 0 & 0 \\
    0 & 1 & 0 & 0 & 0 \\
    0 & 0 & 0 & 0 & 0 \\
\end{pmatrix},
&
M_4 &:=\begin{pmatrix}
    0 & 0 & 0 & 0 & 0 \\
    0 & 0 & 0 & 0 & 0 \\
    0 & 0 & 0 & 0 & -1 \\
    0 & 0 & 0 & 0 & 0 \\
    0 & 0 & 1 & 0 & 0 \\
\end{pmatrix}, \\
M_5 &:=\begin{pmatrix}
    0 & 0 & 0 & 0 & 0 \\
    0 & 0 & 0 & 0 & -1 \\
    0 & 0 & 0 & -1 & 0 \\
    0 & 0 & 1 & 0 & 0 \\
    0 & 1 & 0 & 0 & 0 \\
\end{pmatrix}, & S&:=\begin{pmatrix}
    0 \\
    \frac{1}{h^3} m_1 (m_1 \partial_1 h + m_2 \partial_2 h) \\
    \frac{1}{h^3} m_2 (m_1 \partial_1 h + m_2 \partial_2 h)  \\
    0 \\
    0 \\
\end{pmatrix}.
\end{aligned}
\end{equation}
% System \eqref{ShallowWaterMomentum2D} is \emph{not} a symmetric system due to the terms arising from
% \begin{equation}
%     \label{BadTerm}
%    m \Div m.
% \end{equation}
In particular, $ M_1 $ and $ M_2 $ are not symmetric matrices. Therefore, directly applying energy estimate as in Section \ref{1DEnergyEstimates} will yield a loss of regularity. To remedy this, we isolate the non-symmetric part of system \eqref{ShallowWaterMomentum2D} by writing
\begin{equation}
    \label{def:2d-non-sym}
    K:= B_1 \partial_1 m + B_2 \partial_2 m
\end{equation}
with
\begin{equation}
    \label{def:non-sym-matrix-2d}
    B_1 := \begin{pmatrix}
        \frac{2m_1}{h^2} & 0 \\
        \frac{m_2}{h^2} & \frac{m_1}{h^2}
    \end{pmatrix}, \quad B_2 := \begin{pmatrix}
        \frac{m_2}{h^2} & \frac{m_1}{h^2} \\
        0 & \frac{2m_2}{h^2}
    \end{pmatrix}.
\end{equation}
Now, write 
\begin{equation}
    \label{2d-001}
    B_1 =  \underbrace{\begin{pmatrix}
        \frac{2m_1}{h^2} & \frac{m_2}{2h^2} \\
        \frac{m_2}{2h^2} & \frac{m_1}{h^2}
    \end{pmatrix}}_{:= B_\mathrm{1,sym}} + \frac{m_2}{2h^2} J, \qquad B_2 =  \underbrace{\begin{pmatrix}
        \frac{m_2}{h^2} & \frac{m_1}{2h^2} \\
        \frac{m_1}{2h^2} & \frac{2m_2}{h^2}
    \end{pmatrix}}_{:= B_\mathrm{2,sym}} - \frac{m_1}{2h^2}J.
\end{equation}
Here $ J $ is the anti-symmetric matrix given in \eqref{def:anti-matrix}.
Then one can write $ K $ from \eqref{def:2d-non-sym} as
\begin{equation}
    \label{2d-K-good-form}
    K = B_\mathrm{1,sym} \partial_1 m + B_\mathrm{2,sym} \partial_2 m + \underbrace{\frac{m_2}{2h^2} J \partial_1 m - \frac{m_1}{2h^2} J \partial_2 m}_{=:K'}.
\end{equation}
We can rewrite
\begin{equation}
    \label{2d-002}
    \begin{gathered}
    K' = \frac{m_2}{2h^2} J \nabla m_1 - \frac{m_1}{2h^2} J \nabla m_2 + \frac{m_2}{2h^2}(J\partial_1 m - J\nabla m_1) + \frac{m_1}{2h^2} (J\nabla m_2 - J\partial_2 m) \\
    = \frac{m_2}{2h^2} \nabla^\perp m_1 - \frac{m_1}{2h^2} \nabla^\perp m_2 + \frac{m_2}{2h^2} J (\nabla_\mathrm{skw}m)_1 - \frac{m_1}{2h^2} J (\nabla_\mathrm{skw}m)_2,
    \end{gathered}
\end{equation}
where 
\begin{equation}
    \label{2d-003}
    \nabla_\mathrm{skw} m := \nabla m - (\nabla m)^\top = (\partial_1 m - \nabla m_1, \partial_2 m - \nabla m_2) = \begin{pmatrix}
        0 & - \curl m \\
        \curl m & 0
    \end{pmatrix}
\end{equation}
is the anti-symmetric gradient of $ m $, and $(\nabla_\mathrm{skw} m)_i $
is the $ i $-th column of $\nabla_{\mathrm{skw}}m$. Therefore, $ K' $, written in the form of \eqref{2d-002}, is the sum of the vorticity (the $\nabla_\mathrm{skw} m$ terms) and the dual vorticity (the $J\nabla m_j = \nabla^\perp m_j$, $ j = 1,2$ terms). Thus, the shallow water equations in 2D \eqref{ShallowWaterMomentum2D} can be written as
\begin{equation}
    \label{ShallowWaterMomentum2D-1}
    M_0\partial_t U+M_1' \partial_1 U +M_2' \partial_2 U +M_3\partial_1^2 U +M_4\partial_2^2 U +M_5\partial_{1}\partial_2 U=S + S',
\end{equation}
where $ M_3, M_4, M_5, S $ are given as in \eqref{def:M-2d} and
\begin{equation}
    \label{def:M-2d-1}
    \begin{aligned}
    M_0& :=\diag\left(1,\frac{1}{h},\frac{1}{h},1,1\right), \quad    M_1' :=\begin{pmatrix}
        0 & 1 & 0 & 0 & 0 \\
        1 & \frac{2m_1}{h^2} & {\frac{m_2}{2h^2}} & 0 & 0 \\
        0 & {\frac{m_2}{2h^2}} & \frac{m_1}{h^2} & 0 & 0 \\
        0 & 0 & 0 & 0 & 0 \\
        0 & 0 & 0 & 0 & 0 \\
    \end{pmatrix}, \qquad M_2' :=
    \begin{pmatrix}
        0 & 0 & 1 & 0 & 0 \\
        0 & \frac{m_2}{h^2} & {\frac{m_1}{2 h^2}} & 0 & 0 \\
        1 & {\frac{m_1}{2 h^2}} & \frac{2m_2}{h^2} & 0 & 0 \\
        0 & 0 & 0 & 0 & 0 \\
        0 & 0 & 0 & 0 & 0 \\
    \end{pmatrix}, \\
    S' & := \begin{pmatrix}
        0 \\
        - K' \\
        0 \\
        0 \\
    \end{pmatrix} = \begin{pmatrix}
        0 \\
        - \frac{m_2}{2h^2} \nabla^\perp m_1 + \frac{m_1}{2h^2} \nabla^\perp m_2 - \frac{m_2}{2h^2} J (\nabla_\mathrm{skw}m)_1 + \frac{m_1}{2h^2} J (\nabla_\mathrm{skw}m)_2 \\
        0 \\
        0 \\
    \end{pmatrix}.
    \end{aligned}
\end{equation}


%{\color{purple} The energy estimate below is not correct.}

%Now we can define the functionals. Let 
%\begin{equation}
%    \label{2d-def:energy}
%    \begin{aligned}
%    \mathcal F(t) &:= \sum_{\substack{0 \leq i \leq 1\\0\leq i+j \leq 2}}\norm{\partial_t^i \nabla^j h, \frac{\partial_t^i \nabla^j m}{h^{\frac{1}{2}}}, \partial_t^i \nabla^j \nabla h }{L^2}^2, \\
%\mathfrak F(t) & := \norm{h, m, \nabla h}{H^3}^2.
%    \end{aligned}
%\end{equation}

%First, we establish the energy estimate for the 2d shallow water system:
%\begin{lemma}
%    \label{lm:2d-hr} Assume that \eqref{non-degenerate-nD} holds. One has that
%    \begin{equation}
%        \label{lmest-2d-001}
%        \dfrac{d}{dt} \mathcal F(t) \leq \mathcal H( \mathcal F(t), \mathfrak F(t), \norm{\curl m (t)}{H^2}, \norm{\curl \partial_t m (t) }{H^1} )
%    \end{equation}
    % \begin{equation}
    %     \label{propest-2d:}
    %     \sup_{0\leq t \leq T}\lbrace \norm{\partial_t h, \partial_t m, \partial_t \nabla h}{H^1}^2 + \norm{h, m, \nabla h}{H^2}^2 \rbrace \leq C \norm{h_{in}, m_{in}, \nabla h_{in}}{H^3}^2 
    % \end{equation}
    % for some constant $ T, C \in (0,\infty) $. 
%\end{lemma}

%\begin{proof}
    
%\end{proof}

%Next, we establish the vorticity estimate:
%\begin{lemma}
%    \label{lm:2d-vorticity}
%    \begin{equation}
%        \label{lmest-2d-201}
%    \dfrac{d}{dt} \sum_{\substack{0\leq i, j \leq 2 \\ 0 \leq i+j \leq 2}}\norm{\partial_t^i \nabla^j (\frac{\curl u}{h})}{L^2}^2 \leq \mathcal H(F(t), \mathfrak F(t))
%    \end{equation}
%\end{lemma}


%Finally, we have the elliptic estimate:
%\begin{lemma}
%    \label{lm:2d-elliptic}
%    \begin{equation}
%        \label{lmest-2d-301}
%       \norm{\Div m, \nabla^2 h}{H^2} \leq \mathcal H(\mathcal F(t)), 
%    \end{equation}
%    and hence
%    \begin{equation}
%        \label{lmest-2d-302}
%        \mathfrak F(t) \leq \mathcal H(\mathcal F(t),\norm{\curl m}{H^2}), 
%    \end{equation}
%\end{lemma}

\smallskip 

Meanwhile, unlike the one dimensional flow, we will need to track the evolution of the vorticity to handle $ S' $ in the following. Applying $\curl$ to \eqref{MoMeNtUm} yields
\begin{equation}
    \label{Curled}
    \partial_t \curl u+u\cdot \nabla \curl u = - \Div u \curl u \overset{\eqref{MaSs}}{=} \frac{1}{h}(\partial_t h + u \cdot \nabla h) \curl u.
\end{equation}
Therefore, dividing \eqref{Curled} with $ h $ and re-arranging the resultant equation lead to the specific vorticity formula \cite{BreschDesjardinsMetivier,ShkollerVicol,ChenCialdeaShkollerVicol},
\begin{equation}
    \label{CurledTransport}
    \partial_t \theta +u\cdot \nabla \theta =0,
\end{equation}
where
\begin{equation}
    \label{ThetaVariable}
    \theta:=\frac{\curl u}{h}.
\end{equation}

\smallskip 
In addition, 
consider the energy functional
\begin{equation}
    \label{Energy-2d}
    \mathcal{E}_2(t):=\sum_{\substack{i+ j\leq 3, \\ i,j \geq 0}} \norm{\partial_t^i \nabla^{j}(h-1),\partial_t^i\nabla^j m,\partial_t^i\nabla^{j+1} h}{L^2}^2,
\end{equation}
and the intermediate energy 
\begin{equation}
    \label{intrmd-energy-2d}
    \mathfrak{E}_2(t):=\sum_{\substack{i+ j \leq 3, \\ i, j \geq 0}} (M_0\partial_t^i\nabla^j  U(t),\partial_t^i\nabla^j U(t)).
\end{equation}
By the non-degenerate condition \eqref{non-degenerate-nD}, there exists a constant $C_{\underline{h},\overline{h}}\in (0,\infty)$ such that
       \begin{equation}
       \label{Energy-equiv-2d-cauchy}
            \frac{1}{C_{\underline{h},\overline{h}}}\mathcal{E}_2(t)\leq \mathfrak{E}_2(t)\leq C_{\underline{h},\overline{h}}\mathcal{E}_2(t).
        \end{equation}

\subsection{Vorticity estimate}
\label{subsec:2d-vorticity}

Let $ \partial \in \lbrace \partial_1, \partial_2 \rbrace $ be the spatial derivative.  
For $ i + j \leq 3 $, and $ i \geq 0, \ j \geq 0 $, consider the evolution of $ \theta_{ij}:=\partial_t^i \partial^j \theta $ from \eqref{CurledTransport}, 
\begin{equation}
    \label{eq:dd-curl-2d}
    \partial_t \theta_{ij} + u \cdot \nabla \theta_{ij} = - \sum_{\substack{1\leq a \leq i, \\ 1 \leq b \leq j}} \partial_t^{a} \partial^{b}u \cdot \nabla \partial_t^{i-a}\partial^{j-b} \theta,
\end{equation}
where we have omitted the combinatorial constant on the right hand side. 
Then the standard $ L^2 $ estimate of $ \theta_{ij} $ yields
\begin{equation}
    \label{est:theta-ij}
    \begin{gathered}
    \dfrac{d}{dt}\norm{\theta_{ij}}{L^2}^2 \lesssim \norm{ \Div u}{L^\infty} \norm{\theta_{ij}}{L^2}^2 + \sum_{\substack{1\leq a \leq i, \\ 1 \leq b \leq j}} \norm{\partial_t^{a} \partial^{b}u \cdot \nabla \partial_t^{i-a}\partial^{j-b} \theta}{L^2} \norm{\theta_{ij}}{L^2}\\
    \lesssim \norm{u}{H^3} \norm{\theta_{ij}}{L^2}^2 + \sum_{\substack{i+j\leq 3, \\ i,j \geq 0}}\norm{\partial_t^i\partial^j u}{L^2} \sum_{\substack{i+j\leq 3, \\ i,j \geq 0}} \norm{\partial_t^i \partial^j \theta}{L^2}^2.
    \end{gathered}
\end{equation}
Applying Gr\"onwall's inequality the yields
\begin{equation}
    \label{est:theta-2}
\sum_{\substack{i+j\leq 3, \\ i,j \geq 0}} \norm{\partial_t^i \nabla^j \theta (t)}{L^2}  \lesssim \sum_{\substack{i+j\leq 3, \\ i,j \geq 0}} \norm{\partial_t^i \nabla^j \frac{\curl u (0)}{h(0)}}{L^2} \exp (t \sup_{0\leq s \leq t}  \sum_{\substack{i+j\leq 3, \\ i,j \geq 0}}\norm{\partial_t^i\nabla^j u(s)}{L^2} ).
\end{equation}

% As a consequence, given that $u$ is regular enough, we may take $t\in (0,T)$ sufficiently small such that
% \begin{equation}
% \label{theta-bound}
% \norm{\theta}{H^s}\leq 2\norm{\theta_0}{H^s}.
% \end{equation}
Meanwhile, since, $ u = \frac{m}{h} $ and, with direct calculation, 
\begin{equation}
\label{eq:theta-2}
\theta=\frac{\curl u}{h}= \frac{\curl m}{h^2}-\frac{m\cdot \nabla^{\perp}h}{h},
\end{equation}
the following lemma follows with a straightforward calculation:
\begin{lemma}[Vorticity estimate]
\label{lm:vorticity}
For a solution of \eqref{ShallowWaterMomentum2D} (or equivalently, of \eqref{ShallowWaterMomentum2D-1}) that satisfies \eqref{non-degenerate-nD}, one has that, $ \forall \ t\in (0,\infty) $,  
\begin{equation}
    \label{est:vorticity}
    \sum_{\substack{i+j\leq 3, \\ i,j \geq 0}} 
    \norm{\partial_t^i \nabla^j \curl m (t)}{L^2} 
    \lesssim \sup_{0\leq s \leq t} \mathcal P(\mathcal E_2(s)) + C_\mathrm{skw} 
    \exp(t \sup_{0\leq s \leq t} \mathcal P(\mathcal E_2(s))),
\end{equation}
where $ C_\mathrm{skw}\in(0,\infty) $ is a constant depending on 
\begin{equation}
    \label{2d-initial-bound}
    \sum_{\substack{i+j\leq 3,\\i,j \geq 0}} \norm{\partial_t^i \nabla^j \curl u(t=0),\partial_t^i \nabla^j u(t=0) , \partial_t^i \nabla^j h(t=0),\partial_t^i \nabla^{j+1} h(t=0)}{L^2} < \infty. 
\end{equation}
% \begin{equation}
% \label{initial-data-bound}
%      \norm{\curl u(t=0)}{H^3}+ \norm{u(t=0),h(t=0),\nabla h (t=0)}{H^5} \leq \norm{u(t=0),h(t=0),\nabla h (t=0)}{H^5} < \infty. 
% \end{equation}
\end{lemma}


%Now we write \eqref{Momentum2D} as
%\[
%\frac{1}{h} \partial_t m+\frac{m}{h^2}\cdot\nabla m+ \frac{m}{h^2}\Div m+\nabla h-\nabla\Delta h=\frac{m\cdot \nabla h}{h^3}m.
%\]
%We dissect the momentum term $\frac{m}{h^2}\cdot\nabla m+ \frac{m}{h^2}\Div m$ into symmetric and skew-symmetric parts:
%\begin{align*}
%\frac{m}{h^2}\cdot\nabla m+ \frac{m}{h^2}\Div m &= \begin{pmatrix}
%    \frac{m_x }{h^2} & \frac{1}{2} \frac{m_y }{h^2} \\
%    \frac{1}{2} \frac{m_y }{h^2} & 0 \\
%\end{pmatrix} \partial_x m+\begin{pmatrix}
%    0 & \frac{1}{2} \frac{m_x }{h^2} \\
%    \frac{1}{2}\frac{m_x }{h^2} & \frac{m_y }{h^2} \\
%\end{pmatrix} \partial_y m  \\
%&+\begin{pmatrix}
%    0 & -\frac{1}{2} \frac{m_y }{h^2} \\
%    \frac{1}{2} \frac{m_y }{h^2} & 0 \\
%\end{pmatrix} \partial_x m+ \begin{pmatrix}
%    0 & \frac{1}{2} \frac{m_x }{h^2} \\
%    -\frac{1}{2} \frac{m_x }{h^2} & 0 \\
%\end{pmatrix} \partial_y m.
%\end{align*}
%The symmetric part has the correct structure. Observe that we can %write the skew-symmetric part as
%\begin{align}
%\label{Skew}
%&\begin{pmatrix}
%    0 & -\frac{1}{2} \frac{m_y }{h^2} \\
%    \frac{1}{2} \frac{m_y }{h^2} & 0 \\
%\end{pmatrix} \partial_x m+ \begin{pmatrix}
%    0 & \frac{1}{2} \frac{m_x }{h^2} \\
%    -\frac{1}{2} \frac{m_x }{h^2} & 0 \\
%\end{pmatrix} \partial_y m =\frac{1}{2h^2}\bigg[-m^{\perp}(\curl m) -m_x\nabla^{\perp} m_y+m_y\nabla^{\perp} m_x\bigg].
%\end{align}

\subsection{Energy estimate}
\label{subsec:2d-energy-est}

Now we outline the \emph{a priori} energy estimates for \eqref{ShallowWaterMomentum2D-1}. For simplicity we stay in the original variables $U = (h-1,m,\nabla h)^\top.$ 
The left-hand-side of \eqref{ShallowWaterMomentum2D-1} has the correct symmetric structure, and the first source term $S$ is of lower order. Hence, it remains to check the energy estimates of the source term $S'$. Heuristically, this is done by making the following crucial observation: Thanks to \eqref{2d-003}, $ \nabla_\mathrm{skw} m $ can be controlled by the vorticity estimate \eqref{est:vorticity}, and, with respect to the $ L^2 $ inner product, $ \nabla^\perp $ is the duality operator of $ \curl $, i.e., 
\begin{equation}
\label{duality-formula}
    \int \nabla^\perp f \cdot g \,d\vec x = - \int f \curl g \,d\vec x
\end{equation}
for $ f $ and $ g $ being scaler and vector valued functions, respectively. \eqref{duality-formula} is going to help us to close the energy estimates.
\smallskip


% We are now ready for the \emph{a priori} estimates for the 2D flow \eqref{ShallowWaterMomentum2D-1}. This is done via an argument similar to Lemma \ref{LemmaPhi}.
% \begin{lemma}
%     Let $\Phi:\mathbb{R}\times\mathbb{R}^2 \to\mathbb{R}^3$ be the solution to the linear equation
%     \begin{equation}
%         \label{PhiEq-2d}
%         M_0\partial_t \Phi+M_1'\partial_1 \Phi+M_2'\partial_2 \Phi+M_3 \partial_1^2 \Phi+M_4 \partial_2^2 \Phi+M_5 \partial_{12} \Phi =\mathcal{G},
%     \end{equation}
%     where $\mathcal{G}\in L_t^{\infty} L_x^2$. Then for a.e. $t$, we have
%     \begin{equation}
%         \label{LemmaPhiEstimate-2d}
%         \frac{d}{dt} (M_0\Phi,\Phi)\leq C_{\underline{h},\overline{h}} \mathcal{E}_2(t) \norm{\Phi}{L^2}^2 +2\norm{\mathcal{G}}{L^2}\norm{\Phi}{L^2}.
%     \end{equation}
%     \begin{proof}
%          The proof of Lemma \ref{PhiEq-2d} is similar to that of Lemma \ref{PhiEq}. For the sake of completeness, we write down the whole argument. Take the $L^2$-inner product of \eqref{PhiEq-2d} with $2\Phi$ and integrating the resultant by parts leads to 
%          \begin{equation}
%          \label{phi-2d:001}
%              \frac{d}{dt}(M_0\Phi,\Phi)=(\partial_t M_0\Phi,\Phi)+(\partial_1 M_1' \Phi,\Phi)+(\partial_2M_2'\Phi,\Phi)+( M_5\partial_1 \Phi,\partial_2 \Phi)+2(\mathcal{G},\Phi),
%          \end{equation}
%           where we used the fact that $ M_0, M_1',M_2' $ are symmetric matrices and $ M_3,M_4,M_5 $ are constant skew-symmetric matrices. Moreover, observe that
%     \begin{equation*}
%     \begin{aligned}
%         (M_5\partial_1\Phi,\partial_2\Phi)&=\int\partial_2 \Phi\cdot (0,-\partial_1\Phi_5,-\partial_1\Phi_4,\partial_1\Phi_3,\partial_1\Phi_2) d\vec{x} \\
%         &=\int -\partial_1\Phi_5\partial_2 \Phi_2-\partial_1\Phi_4 \partial_2\Phi_3+\partial_1\Phi_3\partial_2\Phi_3+\partial_1\Phi_2\partial_2\Phi_5 d\vec{x} \\
%         & =\int-\partial_1\Phi_2\partial_2\Phi_5 -\partial_1\Phi_3\partial_2\Phi_4 +\partial_1\Phi_3\partial_2\Phi_3+\partial_1\Phi_2\partial_2\Phi_5 d\vec{x} =0,
%         \end{aligned}
%     \end{equation*}
%     through integration by parts. Thus, applying the Hölder inequality on the right-hand side of \eqref{phi-2d:001} yields
%     \begin{equation}
%         \label{phi-2d:002}
%         \frac{d}{dt}(M_0\Phi,\Phi)\leq \norm{\partial_t M_0,\partial_1 M_1,\partial_2 M_2}{L^2} \norm{\Phi}{L^2}^2+2\norm{\mathcal{G}}{L^2} \norm{\Phi}{L^2}.
%     \end{equation}
%     From \eqref{def:M-2d-1} and \eqref{PhiEq-2d}, one has
%     \begin{equation}
%     \begin{aligned}
%        & \norm{\partial_t M_0,\partial_1 M_1',\partial_2 M_2'}{L^{\infty}} \leq C_{\underline{h}}\norm{\partial_t h,\partial_1 m_1,\partial_1m_2,\partial_2m_1,\partial_2m_2,m_1\partial_1h,m_2\partial_1 h,m_1\partial_2 h,m_2\partial_2 h}{L^{\infty}} \\
%        & \leq C_{\underline{h}} \norm{\partial_1 m_1,\partial_1 m_2,\partial_2 m_1,\partial_2 m_2,\partial_1 h,\partial_2 h,m_1,m_2}{L^{\infty}}^2 \leq C_{\underline{h},\overline{h}} \sum_{0\leq \vert\alpha\vert \leq 3} \norm{\partial^{\alpha}(h-1),\partial^{\alpha}m}{L^2}^2 \leq C_{\underline{h},\overline{h}} \mathcal{E}_2(t),
%     \end{aligned}
%     \end{equation}
%     thanks to the Sobolev embedding inequality $\norm{\cdot}{L^{\infty}} \lesssim \norm{\cdot }{H^2}$ in dimension two.
%     \end{proof}
% \end{lemma}

We will only present the high order estimate. From \eqref{ShallowWaterMomentum2D-1}, one can write down
the systems for $U_{ij}:=\partial_t^i\partial^j U$, $ i + j = 3 $, $ i,j \geq 0 $,
% , $U_{21}:=\partial_t^2 \partial^{\beta-2} U$, $U_{12}:=\partial_t\partial^{\beta-1} U$, and $U_{03}:=\partial^{\beta} U$:
\begin{equation}
\label{sys:dd-2d}
% \begin{aligned}
  \underbrace{M_0\partial_t U_{ij} +M_1'\partial_1 U_{ij}+M_2'\partial_2 U_{ij}+M_3\partial_1^2 U_{ij} +M_4 \partial_2^2 U_{ij}+M_5\partial_1\partial_2 U_{ij} =\mathcal{G}_{ij}}_{\text{symmetric part}}+\underbrace{\mathcal{G}'_{ij}}_{\mathclap{\text{non-symmetric part}}}, \\
  % & \text{for} \quad a\in \{ttt,ttz,tzz,zzz\},
    % \end{aligned}
\end{equation}
where 
\begin{equation}
    \begin{aligned}
      &  \mathcal{G}_{30}:= \partial_t^3 S-3\partial_t M_0\partial_t^3U-3\partial_t^2 M_0\partial_t^2 U-\partial_t^3 M_0\partial_tU-3\partial_tM_1\partial_t^2\partial_1 U-3\partial_t^2 M_1\partial_t \partial_1 U\\
      & -\partial_t^3 M_1 \partial_1 U-3\partial_t M_2\partial_t^2 \partial_2 U-3\partial_t^2 M_2\partial_t \partial_2U-\partial_t^3 M_2\partial_2U,\\
      & \mathcal{G}_{21}:= \partial_{t}^2 \partial  S-\partial M_0\partial_t^3 U-2\partial \partial_t M_0\partial_t^2 U-\partial \partial_t^2 M_0\partial_t U-\partial_t^2 M_0\partial_t\partial  U-2\partial_tM_0\partial_t^2 \partial  U \\
      &-\partial M_1\partial_t^2\partial_1 U-2\partial \partial_t M_1\partial_t\partial_1 U -\partial \partial_t^2M_1\partial_1 U-\partial_t^2 M_1\partial_1\partial  U-2\partial_tM_1\partial_t\partial_1\partial  U \\
      &-\partial  M_2\partial_t^2\partial_2 U-2\partial \partial_t M_2 \partial_t\partial_2 U- \partial  \partial_t^2 M_2\partial_2 U-\partial_t^2 M_2\partial_2\partial U-2\partial_tM_2\partial_t\partial_2\partial  U,\\
      & \mathcal{G}_{12}:= \partial_t \partial^2 S-\partial M_0\partial_t^2 \partial  U-2\partial \partial_t M_0\partial_t\partial  U-\partial^2 \partial_t M_0 \partial  U-\partial_t M_0\partial^2\partial_t U \\
      & -2\partial M_0\partial \partial_t^2 U-\partial M_1\partial_t\partial \partial_1 U-2\partial_t\partial  M_1\partial_1\partial  U-\partial^2 \partial_t M_1\partial_1 U-\partial_tM_1\partial^2\partial_1 U\\
      &-2\partial  M_1\partial \partial_t\partial_1 U-\partial M_2\partial_t\partial \partial_2 U-2\partial_t\partial  M_2\partial_2\partial  U-\partial^2 \partial_t M_2\partial_2 U\\
      &-\partial_tM_2\partial^2\partial_2 U-2\partial  M_2\partial \partial_t\partial_2 U, \\
      & \mathcal{G}_{03}:= \partial^3 S-3\partial  M_0\partial^2 \partial_t U-3\partial^2 M_0\partial \partial_t U -\partial^3 M_0\partial_t U-3\partial  M_1\partial^2\partial_1 U \\
      &-3\partial^2 M_1 \partial  \partial_1 U -\partial^3 M_1 \partial_1 U-3\partial  M_2\partial^2 \partial_2 U-3\partial^2 M_2\partial  \partial_2 U-\partial^3 M_2\partial_2 U,\\
    \end{aligned}
\end{equation}
and 
\begin{equation}
    \mathcal{G}'_{30}:=\partial_t^3 S', \quad \mathcal{G}_{21}':=\partial_t^2 \partial S', \quad \mathcal{G}'_{12}:=\partial_t \partial^2 S', \quad \mathcal{G}'_{03}:=\partial^3 S'.
\end{equation}
Here we have written \eqref{sys:dd-2d} into the symmetric part and the non-symmetric part, where the symmetric part can be handled similar to \eqref{sys:dd-V}--\eqref{est:nonlinear-006}. To be more precise, taking the $ L^2 $-inner product of \eqref{sys:dd-2d} with $ U_{ij} $, and applying integration by parts, H\"older's inequality, and Sobolev embedding inequality, lead to 
\begin{equation}
\label{est:2d-U-ij}
    \frac{1}{2}\dfrac{d}{dt}(M_0 U_{ij}, U_{ij}) \lesssim \mathcal P(\mathcal E_2(t)^{\frac{1}{2}}) \norm{U_{ij}}{L^2}^2 + (\mathcal G_{ij}', U_{ij}). 
\end{equation}

% Applying Lemma \ref{PhiEq-2d} to \eqref{ShallowWaterMomentum2D-1}, we have
% \begin{equation}
%     \frac{d}{dt}(M_0 U,U)\leq C_{\underline{h},\overline{h}} \mathcal{E}_2(t) \norm{U}{L^2}^2+\norm{S}{L^2}\norm{U}{L^2}+(S',U).
% \end{equation}
% One can estimate the symmetry contributions similarly to . 


% Computing the estimates for the symmetric part
%Applying \eqref{LemmaPhiEstimate-2d} to \eqref{sys:dd-2d} leads to 
%\begin{equation}
%\label{h2:001-2d}
%\begin{aligned}
%   \frac{d}{dt}\sum_{a\in \{ttt,ttz,tzz,zzz\}} (M_0U_a,U_a)&\leq C_{\underline{h},\overline{h}} \mathcal{E}_2(t)\norm{U_{ttt},U_{ttz}, U_{tzz}, U_{zzz}}{L^2} \\
%   & +2(\norm{S_{ttt},S_{ttz},S_{tzz},S_{zzz}}{L^2}+\norm{S'_{ttt},S'_{ttz},S'_{tzz},S'_{zzz}}{L^2}) \norm{U}{L^2}.
%    \end{aligned}
%\end{equation}
%A straightforward calculation of the lower order terms gives
%\begin{equation}
%\label{h-1-eq-2d}
%\begin{aligned}
%   & \frac{d}{dt}\lbrace(M_0 \partial_t^2U,\partial_t^2 U)+(M_0\partial_t \partial_z U,\partial_t\partial_z U)+(M_0\partial_z^2 U,\partial_z^2 U)\rbrace \\
%   &=2\lbrace(M_0\partial_t^3 U,\partial_t^2U)+(M_0\partial_t^2 \partial_z U,\partial_t \partial_z U)
%   +(M_0\partial_t\partial_z^2 U, \partial_z^2 U)\rbrace \\
%   &+(\partial_t M_0\partial_t^2 U,\partial_t^2 U)+(\partial_t M_0\partial_t \partial_zU,\partial_t\partial_z U)+(\partial_t M_0\partial_z^2 U,\partial_z^2 U)\\
%   & \lesssim \norm{M_0,\partial_t M_0}{L^{\infty}} \norm{\partial_t^2 U,\partial_t^3 U,\partial_t^2 \partial_z U,\partial_t\partial_z U,\partial_t\partial_z^2 U,\partial_z^2 U}{L^2}^2 \\
%   & \leq C_{\underline{h},\overline{h}} \mathcal{E}_2(t) \norm{\partial_t^2 U,\partial_t^3 U,\partial_t^2 \partial_z U,\partial_t\partial_z U,\partial_t\partial_z^2 U,\partial_z^2 U}{L^2}^2.
%    \end{aligned}
%\end{equation}
%From \eqref{phi-2d:002}, \eqref{h2:001-2d} and \eqref{h-1-eq-2d}, we have
%\begin{equation}
%    \frac{d}{dt} \mathfrak{E}_2(t)\leq C_{\underline{h},\overline{h}} (\mathcal{E}(t))^2 +2(\norm{S,\mathcal{G}_{ttt},\mathcal{G}_{ttz},\mathcal{G}_{tzz},\mathcal{G}_{zzz}}{L^2}+\norm{\mathcal{G}_{ttt}',\mathcal{G}_{ttz}',\mathcal{G}_{tzz}',\mathcal{G}_{zzz}'}{L^2})(\mathcal{E}(t))^{\frac{1}{2}}.
%\end{equation}
%We still need to estimate  $\norm{S,\mathcal{G}_{ttt},\mathcal{G}_{ttz},\mathcal{G}_{tzz},\mathcal{G}_{zzz}}{L^2}$ and $\norm{\mathcal{G}_{ttt}',\mathcal{G}_{ttz}',\mathcal{G}_{tzz}',\mathcal{G}_{zzz}'}{L^2}.$ The former is easy to estimate, as $S$, $M_0$, $M_1$ and $M_2$ depend on $U$ but do not depend on the derivatives of $U$, and we have the non-degenerate condition \eqref{non-degenerate-nD}. For $\Psi\in \{M_0,M_1,M_2,S\}$, we have
%\begin{equation}
%    \begin{aligned}
%       & \norm{\partial_t^3 \Psi}{L^2}\leq C_{\underline{h},\overline{h}}(\norm{\partial_z \Psi}{L^{\infty}}\norm{\partial_t^3 V}{L^2}+\norm{\partial_t\Psi}{L^{\infty}}\norm{\partial_z^2 \Psi}{L^{\infty}}\norm{\partial_t^2 V}{L^2}) \\
%       &\leq C_{\underline{h},\overline{h}} \mathcal{P}(\norm{V}{L^{\infty}})(\norm{\partial_t^3 V}{L^2}+\norm{\partial_z^2 V}{L^{\infty}}\norm{\partial_t^2 V}{L^2}),\\
%       & \norm{\partial_t^2\partial_z \Psi}{L^2} \\
%       & \leq C_{\underline{h},\overline{h}} (\norm{\partial_z \Psi}{L^{\infty}}\norm{\partial_t^2 \partial_z V}{L^2}+\norm{\partial_z \Psi}{L^{\infty}}\norm{\partial_z\partial_t \Psi}{L^{\infty}}\norm{\partial_t^2 V}{L^{2}}+\norm{\partial_t\Psi}{L^{\infty}}\norm{\partial_z^2\Psi}{L^{\infty}}\norm{\partial_t\partial_z V}{L^2}) \\
%       & \leq C_{\underline{h},\overline{h}} \mathcal{P}(\norm{V}{L^{\infty}})(\norm{\partial_t^2\partial_z V}{L^2}+\norm{\partial_t V}{L^{\infty}}\norm{\partial_t\partial_zV}{L^2}+\norm{\partial_z\partial_t V}{L^{\infty}}\norm{\partial_t^2 V}{L^2}), \\
%       & \norm{\partial_t \partial_z^2 \Psi}{L^2} \\
%       &\leq C_{\underline{h},\overline{h}} (\norm{\partial_z\Psi}{L^{\infty}}\norm{\partial_t\partial_z^2 V}{L^2}+\norm{\partial_z \Psi}{L^{\infty}}\norm{\partial_z^2 \Psi}{L^{\infty}}\norm{\partial_t^2 V}{L^2}+\norm{\partial_t\Psi}{L^{\infty}}\norm{\partial_z\Psi}{L^{\infty}}\norm{\partial_z^2 V}{L^2})\\
%       & \leq C_{\underline{h},\overline{h}}\mathcal{P}(\norm{V}{L^{\infty}}) (\norm{\partial_t\partial_z^2 V}{L^2}+\norm{\partial_z^2 V}{L^{\infty}}\norm{\partial_t^2 V}{L^2}+\norm{\partial_t V}{L^{\infty}}\norm{\partial_z^2 V}{L^2}), \\
%       & \norm{\partial_z^3 \Psi}{L^2}\leq C_{\underline{h},\overline{h}} (\norm{\partial_z \Psi}{L^{\infty}}\norm{\partial_z^3 V}{L^2}+\norm{\partial_z \Psi}{L^{\infty}}\norm{\partial_t \Psi}{L^{\infty}}\norm{\partial_z^2 V}{L^2}) \\
%       & \leq C_{\underline{h},\overline{h}} \mathcal{P}(\norm{V}{L^{\infty}}) (\norm{\partial_z^3 V}{L^2}+\norm{\partial_z V}{L^{\infty}}\norm{\partial_t V}{L^{\infty}}\norm{\partial_z^2 V}{L^2}).
%    \end{aligned}
%\end{equation}
%Here we recall that $\mathcal{P}(\cdot)$ is a non-decreasing non-negative polynomial of its argument that differs from line to line, and we also recall that lower order terms were already bounded in \eqref{est:nonlinearity-001}--\eqref{est:nonlinearity-004}.
% The symmetric terms in \eqref{sys:dd-2d} satisfy estimates similar to those in the 1D case \eqref{h-1-eq}--\eqref{est:nonlinear-006}: We omit the details here. A direct computation yields
% \begin{equation}
% \label{energy-ineq-2d-1}
%     \frac{d}{dt}\mathfrak{E}_2(t)\leq C_{\underline h,\overline{h}} \mathcal{E}_2(t)+\mathcal{P}(\mathcal{E}_2(t)^{\frac{1}{2}})+(\mathcal{G}'_a,U_a),
% \end{equation}
% where $a\in \{ttt,ttz,tzz,zzz\}$. 

It remains to calculate $ (\mathcal G_{ij}', U_{ij}) $. We will again only focus on the high order terms. Let  $\partial_H\in \{\partial_t^3,\partial_t^2 \partial ,\partial_t\partial^2, \partial^3\}$. Then from \eqref{def:M-2d-1}, one can check that the high order integral in $ (\mathcal G_{ij}', U_{ij}) $ reads
\begin{equation}
\begin{aligned}
\label{2d-high-energy}
   & \int \frac{\partial_H m}{2h^2}\cdot \bigg[-m_2\nabla^{\perp}\partial_H m_1+m_1\nabla^{\perp} \partial_H m_2-m_2 \partial_HJ(\nabla_{\mathrm{skw}}m)_1+m_1\partial_H J(\nabla_{\mathrm{skw}}m)_2  \\
   & -\partial_H m_2 \nabla^{\perp}m_1+\partial_H m_1\nabla^{\perp}m_2-\partial_H m_2 J(\nabla_{\mathrm{skw}}m)_1+\partial_H m_1J(\nabla_{\mathrm{skw}}m)_2 \bigg]d\vec{x} \\
   & \overset{\eqref{duality-formula}}{=}\int \frac{\partial_Hm}{2h^2} \cdot \bigg[-m_2 \partial_HJ(\nabla_{\mathrm{skw}}m)_1+m_1\partial_H J(\nabla_{\mathrm{skw}}m)_2-\partial_H m_2 J(\nabla_{\mathrm{skw}}m)_1+\partial_H m_1J(\nabla_{\mathrm{skw}}m)_2 \bigg] d\vec{x} \\
   &+\int \frac{1}{2h^2} \left[ m_2\partial_H m_1\curl\partial_H m-m_1 \partial_H m_2 \curl\partial_H m \right] d\vec{x} + \mathrm{l.o.t} \\
   & \lesssim \mathcal E_2^2 \times (\norm{\curl \partial_t^3 m}{L^2}+\norm{\curl \partial_t^2 m}{H^1}+\norm{\curl \partial_t m}{H^2}+\norm{\curl m}{H^3}) + \mathrm{l.o.t},
\end{aligned}
\end{equation}
thanks to the duality \eqref{duality-formula}. 
% Since the specific vorticity $\theta$ satisfies the transport equation \eqref{CurledTransport}, one sees, using \eqref{eq:theta-2},
% \begin{equation}
% \label{curl-2d-3time}
%    \norm{\curl \partial_t^3 m}{L^2} \lesssim \norm{\partial_t^3 \theta}{L^2}+\mathcal{E}_2(t)=\norm{\partial_t^3 \left(\frac{u}{h}\right)}{L^2}+\mathcal{E}_2(t)\lesssim \mathcal{E}_2(t)^{\frac{1}{2}}+\mathcal{E}_2(t)+\mathcal{E}_2(t)^{\frac{3}{2}}+\mathcal{E}_2(t)^2.
% \end{equation}
Here the lower order integral $ \mathrm{l.o.t} $ can be treated similarly as before. 

\smallskip 

Collecting all estimates above as well as the lower order estimates, using \eqref{est:vorticity} from Lemma \ref{lm:vorticity}, we conclude that
% We conclude that the vorticity estimate \eqref{est:vorticity} allows us to control $\norm{\curl \partial_t^3 m}{H^3}$. A similar computation for other terms yields
% \begin{equation}
% \label{curl-2d-others}
%     \norm{\curl \partial_t^2 m}{H^1}+\norm{\curl \partial_t m}{H^2}+\norm{\curl m}{H^3}\lesssim \mathcal{E}_2(t)^{\frac{1}{2}}+\mathcal{E}_2(t)+\mathcal{E}_2(t)^{\frac{3}{2}}+\mathcal{E}_2(t)^2.
% \end{equation} By the vorticity estimate \eqref{est:vorticity} and the bounds \eqref{curl-2d-3time}, \eqref{curl-2d-others}, we have
% \begin{equation}
%     \begin{aligned}
%         \sum_{a\in \{ttt,ttz,tzz,zzz\}} (\mathcal{G}_a',U_a) & \lesssim \norm{m}{H^3}^2 (\mathcal{E}_2(t)^{\frac{1}{2}}+\mathcal{E}_2(t)+\mathcal{E}_2(t)^{\frac{3}{2}}+\mathcal{E}_2(t)^2) \\
%         & \lesssim \mathcal{E}_2(t) (\mathcal{E}_2(t)^{\frac{1}{2}}+\mathcal{E}_2(t)+\mathcal{E}_2(t)^{\frac{3}{2}}+\mathcal{E}_2(t)^2).
%     \end{aligned}
% \end{equation}
% By the energy equivalence \eqref{Energy-equiv-2d-cauchy}, we end up with
% \begin{equation}
%     \label{energy-ineq-2d-2}
%     \sum_{a\in \{ttt,ttz,tzz,zzz\}} (\mathcal{G}_a',U_a)\lesssim \mathfrak{E}_2(t) (\mathfrak{E}_2(t)^{\frac{1}{2}}+\mathfrak{E}_2(t)+\mathfrak{E}_2(t)^{\frac{3}{2}}+\mathfrak{E}_2(t)^2).
% \end{equation}
% Therefore, there exists some existence time $T\in (0,\infty)$ and a non-negative non-decreasing polynomial $\mathcal{P}_2$ such that 
\begin{equation}
    \label{energy-ineq-2d-3}
    \frac{d}{dt}\mathfrak{E}_2(t)\leq C_{\underline h,\overline{h}} \mathcal P(\mathfrak{E}_2(t)^{\frac{1}{2}}) \bigl\lbrack 1 + C_\mathrm{skw}
    \exp(t \sup_{0\leq s \leq t} \mathcal P(\mathfrak E_2(s)^{\frac{1}{2}})) \bigr\rbrack,
    \end{equation}
where we have used \eqref{Energy-equiv-2d-cauchy}.


Therefore, similarly to \eqref{initial_data_time} and \eqref{2d-initial-bound}, one can check that given sufficient regularity $(h_0-1,m_0)\in H^7(\mathbb{R})\times H^6(\mathbb{R}^2)$, we can close both the temporal and spatial energy estimates in dimension two. Hence, we have shown the following: 
%To estimate $\norm{\mathcal{G}_{ttt}',\mathcal{G}_{ttz}',\mathcal{G}_{tzz}',\mathcal{G}_{zzz}'}{L^2}.$, we use \eqref{2d-high-energy} and \eqref{est:vorticity} to deduce 
%\begin{equation}
%\begin{aligned}
%   & \norm{\mathcal{G}_{ttt}'}{L^2}\leq C_{\underline{h},\overline{h}}(\norm{\curl m}{L^{\infty}} \norm{\partial_t^3 m}{L^2}+\norm{\partial_t^2 U}{L^2}\norm{\partial_t \curl m}{L^{\infty}}+\norm{\partial_t U}{L^{\infty}}\norm{\partial_t^2 \curl m}{L^2}), \\
%   & \norm{\mathcal{G}_{ttz}'}{L^2}\leq C_{\underline{h},\overline{h}}( \norm{\curl m}{L^{\infty}}\norm{\partial_t^2 \partial_z m}{L^2}^2+\norm{\partial_t\partial_z U}{L^2}\norm{\partial_t \curl m}{L^{\infty}}+\norm{\partial_t^2 U}{L^2}\norm{\partial_z \curl m}{L^{\infty}}\\
%   & +\norm{\partial_t U}{L^{\infty}}\norm{\partial_t\partial_z\curl m}{L^2}+\norm{\partial_z U}{L^{\infty}}\norm{\partial_t^2 \curl m}{L^2}), \\
%   &\norm{\mathcal{G}_{tzz}'}{L^2}\leq C_{\underline{h},\overline{h}} (\norm{\curl m}{L^{\infty}}\norm{\partial_t^2 \partial_z m}{L^2}+\norm{\partial_t\partial_z U}{L^2}\norm{\partial_z\curl m}{L^{\infty}}+\norm{\partial_z^2 U}{L^{\infty}}\norm{\partial_t \curl m}{L^{\infty}})\\
%   &+\norm{\partial_t U}{L^{\infty}}\norm{\partial_z^2\curl m}{L^2}+\norm{\partial_z U}{L^{\infty}}\norm{\partial_t\partial_z \curl m}{L^2}, \\
%   &\norm{\mathcal{G}'_{zzz}}{L^2}\leq C_{\underline{h},\overline{h}} (\norm{\curl m}{L^{\infty}}\norm{\partial_z^3 m}{L^2}^2+\norm{\partial_z^2 U}{L^2}\norm{\partial_z \curl m}{L^{\infty}}+\norm{\partial_z U}{L^{\infty}}\norm{\partial_z^2 \curl m}{L^2}).\\
%\end{aligned}
%\end{equation}
%Hence, using the Sobolev embedding inequality, we conclude with
%\begin{equation}
%\label{est:nonlinear-2d}
%\begin{aligned}
%&\norm{S,\mathcal{G}_{ttt},\mathcal{G}_{ttz},\mathcal{G}_{tzz},\mathcal{G}_{zzz}}{L^2}+\norm{\mathcal{G}_{ttt}',\mathcal{G}_{ttz}',\mathcal{G}_{tzz}',\mathcal{G}_{zzz}'}{L^2} \\
%&\leq C_{\underline{h},\overline{h}} \bigg[\sum_{\Psi\in \{M_0,M_1,M_2,S\}}\lbrace\norm{\partial_t\Psi,\partial_z \Psi}{L^{\infty}}\norm{\partial_t^3 V,\partial_z^3 V,\partial_t^2 V\partial_z V,\partial_t V\partial_z^2 V}{L^2}+\norm{\partial_t\partial_z \Psi,\partial_z^2 \Psi}{L^{2}}\norm{\partial_t V,\partial_z V}{L^{\infty}}\rbrace\\
%& +\norm{S,\partial_t S,\partial_t^2 S,\partial_t^3 S,\partial_zS,\partial_z^2S,\partial_z^3 S,\partial_t\partial_z S,\partial_t\partial_z^2 S,\partial_t^2 \partial_z S}{L^2}\\
%& +\norm{\partial_t^2 \curl m,\partial_t\partial_z \curl m,\partial_z^2 \curl m}{L^2}\norm{\partial_tU,\partial_t^2 U,\partial_t^3 U,\partial_z U,\partial_z^2 U,\partial_z^3 U,\partial_t\partial_z U,\partial_t^2 \partial_z U}{L^2}\\
%& +\sum_{0\leq \vert\alpha\vert\leq 2}\norm{\partial^{\alpha}\curl m,\partial^{\alpha+1}\curl m,\partial^{\alpha}\partial_t\curl m}{L^2}\norm{\partial_tU,\partial_t^2 U,\partial_t^3 U,\partial_z U,\partial_z^2 U,\partial_z^3 U,\partial_t\partial_z U,\partial_t^2 \partial_z U}{L^2}\bigg] \\
%& \leq \mathcal{P}(\mathfrak{E}_2(t))^{\frac{1}{2}}.
%\end{aligned}
%\end{equation}
%In conclusion, from \eqref{h-1-eq-2d} and \eqref{est:nonlinear-2d}, we deduce that there exists $T>0$ such that
%\begin{equation}
%    \frac{d}{dt} \mathfrak{E}_2(t)\leq \mathcal{P}(\mathfrak{E}_2(t)^{\frac{1}{2}}),
%\end{equation}
%for all $t\in (0,T)$. 
\begin{proposition}
    Under the assumption \eqref{non-degenerate-nD}, there exists the solution to \eqref{ShallowWaterMomentum2D} satisfies
    \begin{equation}
        \sup_{0\leq t\leq T} \mathfrak{E}_2(t)\leq 2\mathfrak{E}_2(0),
        \end{equation}
        and, thanks to \eqref{Energy-equiv-2d-cauchy}, also satisfies        \begin{equation}
          \sup_{0\leq t\leq T}  \mathcal{E}_2(t)\leq C_{\underline{h},\overline{h}} \mathcal{E}_2(0),
        \end{equation}
        for some $T\in (0,\infty)$ and some $C_{\underline{h},\overline{h}}\in (0,\infty) $ depending only on the initial data and $\underline{h}$, $\overline{h}$.
\end{proposition}
Meanwhile, \eqref{non-degenerate-nD} can be closed following the arguments as in Proposition \ref{prop:a-priori-assumption}. Thus repeating the vanishing viscosity arguments as in Section \ref{vv}, we have shown the following: 
\begin{proposition}
\label{prop:2d-local-wellposedness}
    Consider initial data $(h_0-1,m_0)\in H^7(\mathbb{R}) \times H^6(\mathbb{R}^2)$. There exists some time $T\in (0,\infty)$ such that there exists a unique local-in-time solution $(h-1,m)\in L^{\infty}([0,T];H^4(\mathbb{R}))\times L^{\infty}([0,T];H^3(\mathbb{R}^2))$ to \eqref{ShallowWaterMomentum2D}, which satisfies
    \begin{equation}
        \sup_{0\leq t\leq T}\sum_{\substack{i+\vert\alpha\vert\leq 3, \\ i,\vert\alpha\vert\geq 0}} \norm{\partial_t^i\partial^{\alpha}(h-1),\partial_t^i\partial^{\alpha} m, \partial_t^i\partial^{\alpha}\nabla h}{L^2}^2 \leq C_{\underline{h},\overline{h}}\mathcal{E}_2(0),
    \end{equation}
    for some constant $C_{\underline{h},\overline{h}}\in (0,\infty)$ depending only on $\underline{h},$ $\overline{h}$. Moreover, following a calculation similar to \eqref{initial_data_time}, we have
    \begin{equation}
        \mathcal{E}_2(0)\lesssim \mathcal{P} (\norm{h_0-1}{H^7},\norm{m_0}{H^6}).
    \end{equation}
\end{proposition}

%For brevity we only showcase how to deal with the troublesome skew-symmetric terms in \eqref{Skew}. In the $(\eta-1,\omega,\phi)$-variable, we need to perform the $H^r$ energy estimates of
%\[
%\frac{1}{2\eta^2}\bigg[-\omega^{\perp}(\curl \omega) -\omega_x\nabla^{\perp} \omega_y+\omega_y\nabla^{\perp} \omega_x\bigg],
%\]
%in both space and time, with $r\geq 3$ being an integer. The justified assumption \eqref{non-degenerate-nD} allows us to ignore the lower order term $\frac{1}{2\eta^2}$. First consider the estimates in space: We apply $\partial^r$ to \eqref{Skew}, dot with $\partial^r \omega$, and integrate the subsequent expression. In dimension two, the terms with the leading order of derivatives in the energy estimates are
%\begin{equation}
%\label{HighestOrderTerms}
%\int[(-\partial^r \omega_y\nabla^{\perp} \omega_x+\partial^r \omega_x \nabla^{\perp}\omega_y) \cdot \partial^r \omega+(-\omega_x\partial^r\nabla^{\perp} \omega_y+\omega_y\partial^r\nabla^{\perp} \omega_x)\cdot \partial^r \omega] d\vec{x}.
%\end{equation}
%The first term $(-\partial^r \omega_y\nabla^{\perp} \omega_x+\partial^r \omega_x \nabla^{\perp}\omega_y) \cdot \partial^r \omega$ can be directly estimated using the Sobolev embedding inequality $\norm{\cdot}{L^{\infty}}\lesssim \norm{\cdot}{H^2}$:
%\[
%\int (-\partial^r \omega_y\nabla^{\perp} \omega_x+\partial^r \omega_x \nabla^{\perp}\omega_y) \cdot \partial^r \omega d\vec{x}\lesssim \Vert\partial^r \omega\Vert_{L^2}^2 \Vert \nabla \omega\Vert_{L^{\infty}}\lesssim \mathcal{E}_2(t)^{\frac{3}{2}}.
%\]
%Integrating the second term $(-\omega_x\partial^r\nabla^{\perp} \omega_y+\omega_y\partial^r\nabla^{\perp} \omega_x)\cdot \partial^r \omega$ by parts gives
%\begin{align*}
%& \int(-\omega_x\partial^r\nabla^{\perp} \omega_y+m_y\partial^r\nabla^{\perp} \omega_x)\cdot \partial^r \omega d\vec{x}=\int \omega\cdot \partial^r \curl \omega\cdot \partial^r \omega^{\perp} d\vec{x} + \ \text{l.o.t} \\
%&\lesssim \norm{\partial^r \curl \omega}{L^2} \norm{\partial^r \omega^{\perp}}{L^2} \norm{\omega}{L^{\infty}} +\norm{\text{l.o.t}}{H^s}\lesssim \mathcal{E}_2(t) \norm{ w_0}{H^s}+\norm{\text{l.o.t}}{H^s},
%\end{align*}
%where
%\[
%\text{l.o.t}=\int \partial^r \omega_x (-\partial_y \omega_x \partial^r \omega_y +\partial_y \omega_y \partial^r \omega_x)+\partial^r \omega_y (\partial_x \omega_x \partial^r \omega_y-\partial_x \omega_y\partial^r \omega_x) d\vec{x}
%\]
%is of lower order, and is thus controlled by $\mathcal{E}_2(t)$. On the other hand, applying $\partial_t^r$ to \eqref{Skew}, multiplying with $\partial_t^r \omega$ and integrating gives
%\begin{equation}
%\label{HighestOrderTime}
%\begin{aligned}
%&\int \bigg[\partial_t^r \omega\cdot [-\partial_t^r \omega^{\perp}(\curl \omega)-\partial_t^r \omega_x\nabla^{\perp}\omega_y+\partial_t^r \omega_y \nabla^{\perp}\omega_x] \\
%&+\partial_t^r \omega\cdot[-\omega^{\perp}\partial_t^r\curl \omega-\omega_x\partial_t^r \nabla^{\perp}\omega_y+\omega_y\partial_t^r \nabla^{\perp}\omega_x] \bigg] d\vec{x},
%\end{aligned}
%\end{equation}
%as the highest order terms. The first term in \eqref{HighestOrderTime} is bounded:
%\begin{align*}
%& \int [\partial_t^r \omega\cdot [-\partial_t^r \omega^{\perp}(\curl \omega)-\partial_t^r \omega_x\nabla^{\perp}\omega_y+\partial_t^r \omega_y \nabla^{\perp}\omega_x] d\vec{x} \lesssim \norm{\partial_t^r \omega}{L^2}^2 (\norm{\curl \omega}{L^{\infty}}+\norm{\nabla \omega}{L^{\infty}}) \\
%&\lesssim \mathcal{E}_2(t)\cdot(2\norm{w_0}{H^s}+\mathcal{E}_2(t)^{\frac{1}{2}}),
%\end{align*}
%and is thus controlled by the energy \eqref{Energy_higher_d}. Integrating the second term in \eqref{HighestOrderTime} by parts yields
%\begin{align*}
%    &\int (\partial_t^r \omega\cdot[-\omega^{\perp}\partial_t^r\curl \omega-\omega_x\partial_t^r \nabla^{\perp}\omega_y+\omega_y\partial_t^r \nabla^{\perp}\omega_x)] d\vec{x} \\
%    & = \int [\Div \omega \cdot (\partial_t^r m)^2-(\partial_x \omega_y+\partial_y \omega_x)\partial_t^r \omega_x \partial_t^r \omega_y] d\vec{x}\lesssim \norm{\partial_t^r \omega}{L^2}^2 \norm{\Div \omega}{L^{\infty}}\lesssim \mathcal{E}_2(t)^{\frac{3}{2}}.
%\end{align*}
%All other coefficients in \eqref{ShallowWaterMomentum2D} have the same (correct) structure as their counterparts in the 1D case, so they can be directly controlled in the energy estimates. \\
\section{Remark on the regularity \eqref{n-d-initial-regularity} for 2-dimensional flows}
\label{Regularity-remark}
The purpose of Section \ref{Regularity-remark} is to demonstrate the optimal regularity argument in Section \ref{Sect:WP} fails in dimension two. Namely, we cannot relax the regularity of the initial data from $(h_0 -1, u_0) \in H^7(\mathbb{R})\times H^6(\mathbb{R}^2)$ to $(h_0 -1, u_0) \in H^4(\mathbb{R})\times H^3(\mathbb{R}^2)$. The integer Sobolev index corresponding to $n=2$ is $s=3$: The corresponding evolution energy and elliptic energy are
\begin{equation}
    \label{EnergyEvolution-Odd}
    \mathcal{F}_{\text{bad}} (t):=\sum_{\substack{0\leq i\leq 1, \\ 0\leq i+\vert\alpha\vert\leq 2}}\norm{\partial_t^i\partial^{\alpha}(h-1),\frac{\partial_t^i \partial^{\alpha} m}{\sqrt{h}},\partial_t^i \partial^{\alpha}  \nabla h}{L^2}^2,
\end{equation}
and 
\begin{equation}
    \label{EnergyElliptic-Odd}
    \mathfrak{F}_{\text{bad}}(t):=\sum_{0\leq \vert \alpha \vert \leq 3} \norm{\partial^{\alpha}(h-1),\frac{\partial^{\alpha}m}{\sqrt{h}}, \partial^{\alpha}\nabla h}{L^2}^2,
\end{equation}
respectively. To integrate by parts in the energy estimates, one needs to compute, for $\vert \alpha \vert=1$,
\begin{equation}
    \partial_t\partial^{\alpha} \left(\frac{1}{h}\partial_t m\right) \partial_t\partial^{\alpha} m=[\underbrace{\partial^{\alpha} \left(\frac{1}{h}\right)}_{\in L^{\infty}} \boxed{\partial_t^2 m} ]\underbrace{\partial_t\partial^{\alpha} m}_{\in L^2} +\cdots .%+[\underbrace{\partial_{t}\partial^{\alpha} \left(\frac{1}{h}\right)}_{\in L^{\infty}}\underbrace{\partial_t m}_{\in L^2} ]\underbrace{\partial_{t}\partial^{\alpha} m}_{\in L^2} +\underbrace{\partial_t \left(\frac{1}{h}\right)}_{\in L^{\infty}} \underbrace{(\partial_t\partial^{\alpha} m)^2}_{\in (L^2)^2}.
\end{equation}
We require $\partial_t^2m\in L_x^2$: A direct calculation using \eqref{ShallowWaterMomentumMomentum} gives
\begin{equation}
\label{2-time-diff-momentum}
\begin{aligned}
   & \frac{1}{h}{\partial_t^2 m}+\cdots %-\frac{1}{h^2}\partial_tm\partial_t h+\frac{\partial_t m}{h^2}\cdot \nabla m+\frac{m}{h^2}\cdot \partial_t \nabla m-\frac{2m\partial_t h}{h^3}\cdot \nabla m+\frac{\partial_t m}{h^2}\cdot \Div m +\frac{m}{h^2}\cdot \partial_t \Div m \\
   ={\partial_t \nabla \Delta h}+\cdots . %=\frac{2(\partial_t m\cdot \nabla )h}{h^3} m+\frac{(m\cdot \nabla)\partial_t h}{h^3}m-\frac{3(m\cdot \nabla)h}{h^4}m\partial_t h+\frac{2m\partial_t h}{h^3} \cdot \Div m- \partial_t \nabla h. 
\end{aligned}
\end{equation}
However, the term $\partial_t \nabla\Delta h$ is not in the evolution energy \eqref{EnergyEvolution-Odd}. Consequently, $\mathcal{F}_{\text{bad}}$ and $\mathfrak{F}_{\text{bad}}$ cannot satisfy an inequality of the form \eqref{EllipticEstimate1} due to a loss of regularity. Hence, we conclude that it is not possible to loosen the regularity \eqref{n-d-initial-regularity} to the optimal regularity $(h_0-1,m_0)\in H^4(\mathbb{R})\times H^3(\mathbb{R}^2)$ without imposing additional regularity assumptions on $m_0$.
\begin{remark}
    We do not try to optimize the regularity as in Section \ref{Sect:WP} for dimensional higher than $ 1 $ hereafter.  
\end{remark}
\section{$ n $-dimensional irrotational flows, $ n \geq 3 $}
\label{sec:irrotational-flow}

For flows in dimensions $ 3 $ or larger, one can still write down the system \eqref{ShallowWaterEquations} as a sum of the symmetric and non-symmetric parts as in \eqref{ShallowWaterMomentum2D-1}. However, the corresponding vorticity does not satisfy a transport equation similar to \eqref{CurledTransport}. 

To be more precise,  let $ J_{ij}, I_{ij} \in \mathbb R^{n\times n} $, $ 1\leq i,j \leq n $, be matrices defined by
\begin{equation}
    \label{def:j-ij}
    (J_{ij})_{lk} := - \delta_{li}\delta_{kj} +  \delta_{ki}\delta_{lj},
    % \begin{cases}
    %     1 & l = i, k = j, \\
    %     -1 & k = i, l = j, \\
    %     0 & \text{otherwise},
    % \end{cases} 
    \qquad (I_{ij})_{lk} := \delta_{li}\delta_{kj} +  \delta_{ki}\delta_{lj}.
    % \begin{cases}
    %     1 & l = i, k = j, \\
    %     1 & k = i, l = j, \\
    %     0 & \text{otherwise}
    % \end{cases}
\end{equation}


\smallskip 

In analogy to \eqref{def:2d-non-sym}, we focus on the term 
\begin{equation}
    \label{def:n-d-non-symmetry}
    K:= \frac{m}{h^2} \cdot \nabla m + \frac{m}{h^2} \Div m =  \sum_{i=1}^n B_i \partial_i m,
\end{equation}
where 
\begin{equation}
    \label{def:B-i}
    (B_i)_{lk} = \frac{m_i}{h^2} \delta_{lk}  + 
        \frac{m_l}{h^2} \delta_{ki}.
\end{equation}
Then one can write
\begin{equation}
    \label{n-d-001}
    B_i = B_\mathrm{i,sym} + \sum_{j=1}^n A_{ij},
\end{equation}
where
\begin{equation}
    \label{n-d-002}
    \begin{aligned}
        B_\mathrm{i,sym}&:= \frac{m_i}{h^2} \mathbb I_n +
        \sum_{j=1}^n\frac{m_j}{2 h^2} I_{ij}, \\
        A_{ij} & :=
        \frac{m_j}{2 h^2} J_{ij}. 
    \end{aligned}
\end{equation}
Indeed,
\begin{align*}
    ( I_{ij} + J_{ij} )_{lk} & = 2 \delta_{ki} \delta_{lj}, \\
    \intertext{by definition, and therefore}
    (B_\mathrm{i,sym} + \sum_{j=1}^n A_{ij})_{lk} & = \frac{m_i}{h^2} \delta_{lk} + \sum_{j=1}^n \frac{m_j}{2h^2} (I_{ij} + J_{ij})_{lk} \\
    & = \frac{m_i}{h^2} \delta_{lk} + \sum_{j=1}^n \frac{m_j}{2h^2} 2 \delta_{ki}\delta_{lj} 
    = \frac{m_i}{h^2}\delta_{lk} + \frac{m_l}{h^2}\delta_{ki} \overset{\eqref{def:B-i}}{=} (B_i)_{lk}.
\end{align*}
Therefore, substituting \eqref{n-d-001} and \eqref{n-d-002} into \eqref{def:n-d-non-symmetry} yields 
\begin{equation}
    \label{n-d-003}
    \begin{aligned}
    K & = \sum_{i=1}^n B_\mathrm{i,sym} \partial_i m + \sum_{i,j=1}^n \frac{m_j}{2h^2} J_{ij} \partial_i m \\ & = \sum_{i=1}^n B_\mathrm{i,sym} \partial_i m + \sum_{i,j=1}^n \frac{m_j}{2h^2} J_{ij} \underbrace{(\partial_i m - \nabla m_i)}_{= (\nabla_\mathrm{skw} m)_i }  + \sum_{i,j=1}^n \frac{m_j}{2h^2} J_{ij} \nabla m_i,
    \end{aligned}
\end{equation}
where
\begin{equation}
    \label{n-d-005}
    J_{ij} \nabla \cdot m = - \partial_j m_i + \partial_i m_j = (\nabla_\mathrm{skw} m)_{ij},
\end{equation}
is the $ ij $-th component of the anti-symmetric gradient. 
In particular, \eqref{n-d-005} implies the duality
\begin{equation}
    \label{duality-n-d}
    \int J_{ij}\nabla f \cdot g \,d\vec x = - \int f (\nabla_\mathrm{skw} g)_{ij} \,d\vec x,
\end{equation}
in analogy to \eqref{duality-formula}.

Now we are ready to write the system \eqref{ShallowWaterMomentum} into symmetric form in higher dimensions. Let $ U := (h-1, m, \nabla h) = (h-1, m_1, m_2, \cdots, m_n, \partial_1 h, \partial_2 h, \cdots , \partial_n h) \in \mathbb R^{2n+1} $. Write, for $ i,j = 1,2,\cdots , n $,  
\begin{equation}
    \label{def:coefficient}
    \begin{aligned}
        M_0 &:= \diag (1, \frac{1}{h}\mathbb I_n, \mathbb I_n), &
        M_i &:= \begin{pmatrix}
            0 & \vec e_i^\top & 0\\ 
            \vec e_i & B_\mathrm{i,sym} & 0 \\
            0 & 0 & 0 \times \mathbb I_n
        \end{pmatrix}, \\
        M_{ij} &:= \begin{pmatrix}
            0 & 0 & 0 \\
            0 & 0 & M_{ij}'' \\
            0 & M_{ij}' & 0
        \end{pmatrix}, & & % \\
        % \text{with} & \ (M_{ij}')_{lk} := \delta_{li}\delta_{kj} \ \text{for}\ l, k = 1,2,\cdots n, & \text{and}  & 
    \end{aligned}
\end{equation}
with 
\begin{equation}
    \label{def:m-ij'''}
    (M_{ij}')_{lk} := \delta_{li}\delta_{kj} \ \text{for}\ l, k = 1,2,\cdots n, \quad \text{and}  \quad M_{ij}'' = - (M_{ij}')^\top,
\end{equation}
and 
\begin{equation}
    \label{def:source}
    S := \begin{pmatrix}
    0 \\ \frac{(m\cdot\nabla)h}{h^3}  m \\ 0
    \end{pmatrix}, \qquad 
    S'':= \begin{pmatrix}
        0 \\ 
        - \sum_{i,j=1}^n \frac{m_j}{2h^2} J_{ij}  (\nabla_\mathrm{skw} m)_i - \sum_{i,j=1}^n \frac{m_j}{2h^2} J_{ij} \nabla m_i  \\
        0
    \end{pmatrix}.
\end{equation}
It is easy to verify that $ M_i $, $ i = 0,1,2,\cdots n $, are symmetric matrices and $ M_{ij} $, $ i,j = 1,2,\cdots n $ are constant anti-symmetric matrices. Then the system \eqref{ShallowWaterMomentum} can be written as
\begin{equation}
    \label{n-d-006}
    \underbrace{M_0 \partial_t U + \sum_{i=1}^n M_i \partial_i U + \sum_{i,j =1}^n M_{ij} \partial_{ij} U = S}_{\text{symmetric part}} \qquad  + \underbrace{S''}_{\mathclap{\text{non-symmetric part}}}. 
\end{equation}
While the symmetric part can be handled as in the one and two dimensional cases, the non-symmetric part requires the estimates of 
\begin{equation}
\label{id:m-2-u-vorticity}
    \nabla_\mathrm{skw} m = h {\nabla_\mathrm{skw} u} + u (\nabla h)^\top - \nabla h u^\top. 
\end{equation}
Meanwhile, from \eqref{MoMeNtUm}, one can calculate, for $ \forall \ l,k = 1,2,\cdots, n $, 
\begin{gather*}
    \partial_t (\partial_k u_l - \partial_l u_k) = - \partial_k (u\cdot \nabla u_l) + \partial_l (u\cdot \nabla u_k) \\ = - u \cdot \nabla (\partial_k u_l - \partial_l u_k) - \sum_{j=1}^n ( \partial_k u_j \partial_j u_l - \partial_l u_j \partial_j u_k) \\
    = - u \cdot \nabla (\partial_k u_l - \partial_l u_k) - \sum_{j=1}^n \lbrack  (\partial_k u_j - \partial_j u_k) \partial_j u_l + (\partial_j u_l - \partial_l u_j) \partial_j u_k \rbrack.
\end{gather*}
That is,
\begin{equation}
\label{CurledTransport-n-d-before}
    \partial_t \nabla_\mathrm{skw} u + u \cdot \nabla \nabla_\mathrm{skw} u= - \nabla_\mathrm{skw} u (\nabla u)^\top - \nabla u \nabla_\mathrm{skw} u.
\end{equation}
In particular, if $ \nabla_\mathrm{skw} u(t=0)=0 $, one can conclude from \eqref{CurledTransport-n-d-before} that, for all $ \forall \ t \geq 0 $,
\begin{equation}
    \label{irrotation}
    \nabla_\mathrm{skw} u(t) \equiv 0.
\end{equation}
Thus, from \eqref{id:m-2-u-vorticity} and \eqref{irrotation}, one can conclude that 
\begin{equation}
\label{est:skw-m}
    \nabla_\mathrm{skw}m = u (\nabla h)^\top - \nabla h u^\top \simeq U^2,
\end{equation}
where no loss of regularity appears. 
Thus, repeating the arguments as in Section \ref{sec:2d} leads to the following:
\begin{proposition}
    \label{prop:n-d}
    For $ n \geq 3 $, let 
    \begin{equation}
    \label{def:s-for-n-d}
        \frac{n}{2}+1 < s_n:= \begin{cases}
            \frac{n}{2}+ \frac{3}{2} & \text{if $ n $ is odd}, \\
            \frac{n}{2}+2 & \text{if $ n $ is even}.
        \end{cases}
    \end{equation}
    Consider initial data $ (h_0-1, u_0) \in H^{2 s_n+1}(\mathbb R) \times H^{2 s_n}(\mathbb R^n) $ for the system \eqref{ShallowWaterEquations} with $ \nabla_\mathrm{skw} u_0 = 0 $ and \eqref{h_initial-bound}. Then there exists some time $T\in (0,\infty)$ such that there exists a unique local-in-time solution $(h-1,m)\in L^{\infty}([0,T];H^{s_n+1}(\mathbb{R}))\times L^{\infty}([0,T];H^{s_n}(\mathbb{R}^n))$ to \eqref{ShallowWaterEquations}.
\end{proposition}


\section*{Acknowledgments}
Part of this work was compiled during the second author's visit to the Fields Institute during the program \emph{Thematic Program on Shocks and Singularities: Nonlinear Evolution Equations in Physical and Life Sciences}. BY graciously acknowledges the financial support by the Fields Institute, and thanks the organizers for their warm hospitality. 

\section*{Conflict of Interest}
On behalf of all authors, the corresponding author states that there is no conflict of interest.

\section*{Data Availability}
Data sharing is not applicable to this article as no new data were created or analyzed in this study.

\bibliographystyle{plain}
\bibliography{main}







\end{document}
